\subsection{调和综合}

设 A 是一个交换 Banach 代数。回顾：若 M 是 \(\mathsf{X}(\mathsf{A})\) 的一个子集，我们用 \(\Upsilon(\mathsf{M})\) 表示 M 中元素的核的交（参见 I, p.~30）；它是 A 的一个理想。

命题 4. --- 设 A 是一个无根的正则交换 Banach 代数。设 F 是 X(A) 的一个闭子集。使 V(I) = F 的 A 的理想 I 所组成的集合按包含关系排序后有一个最大元，即 Y(F)，还有一个最小元，即由 G(x) 具有与 F 不相交的紧支集的 x ∈ A 组成的集合 J。

按构造，\(\Upsilon(F)\) 是 A 的一个理想，满足 \(V(\Upsilon(F))\) 包含 F，并且它是具有这一性质的 A 的最大理想。由于 F 是闭的，存在 A 的理想 I 使 \(V(I) = F\)；于是有 \(I \subset \Upsilon(F)\)，从而 \(V(\Upsilon(F)) \subset V(I) = F\)，即 \(V(\Upsilon(F)) = F\)。这就证明了第一个断言。

集合 J 是 A 的理想且 V(J) 包含 F。我们来证明 V(J) = F。设 \(\chi \in \mathsf{X}(\mathsf{A})\) 不属于 F。取 \(\chi\) 的一个与 F 不相交的紧邻域 U (TG, I, p.~65, 命题 9 的推论)。依 I, p.~88 命题 1 的断言 (ii)，存在 \(x \in A\) 使 \(\mathcal{G}(x)\) 在 \(\chi\) 处等于 1 而在 U 之外等于 0。于是有 \(x \in J\)，因此 \(\chi \notin \mathrm{V}(\mathrm{J})\)。这表明 V(J) ⊂ F，从而 V(J) = F。

最后，设 I 是 A 的理想且 \(V(I) = F\)。我们证明 \(J \subset I\)。设 \(x \in J\)，并设 C 是 \(\mathcal{G}(x)\) 的支集；集合 C 是 \(X(A)\) 的一个与 F 不相交的紧子集。依命题 3，存在元素 \(u \in I\) 使 \(\mathcal{G}(u) = 1\) 在 C 上成立。于是有 \(\mathcal{G}(x) = \mathcal{G}(ux)\)，又因 A 无根，故 x = ux (I, p.~38, 命题 8)。从而有 \(x \in I\)，这表明 \(J \subset I\)。

推论 1. --- 设 A 是一个无根的正则交换 Banach 代数。设 J 是由 \(x \in \mathrm{A}\) 中使 \(\mathcal{G}(x)\) 具有紧支集的那些元素组成的集合。假设 \(\overline{\mathrm{J}} = \mathrm{A}\)。那么 A 的每个异于 A 的闭理想都含于一个正则极大理想。

若 I 是 A 的闭理想且不含于任何正则极大理想，则 \(V(I) = \varnothing\)，于是 \(I \supset J\)（命题 4 应用于 \(F = \varnothing\)），从而 \(I \supset \overline{J} = A\)。

推论 2. --- 设 A 是一个无根的正则交换 Banach 代数。设 \(x, y \in A\)。若 \(\mathcal{G}(x)\) 的支集是紧的，且含于由特征标 \(\chi\) 组成的集合之中且在其上 \(\mathcal{G}(y)(\chi) \neq 0\)，则 \(x\) 是 \(y\) 在 \(A\) 中的倍元。

设 I 是 A 的理想 Ay。于是 V(I) 是 \(\mathcal{G}(y)\) 的零点集。由于 \(\mathcal{G}(x)\) 的支集 F 是紧的且与 V(I) 不相交，我们有 \(x \in I\)（命题 4 应用于 F）。

定义 2. --- 设 A 是一个交换 Banach 代数。

设 I 是 A 的理想，\(x \in A\)，\(\chi \in \mathsf{X}'(\mathsf{A})\)。我们称 x 在 \(\chi\) 的邻域内属于 I，如果存在元素 \(y \in I\) 使 \(\mathcal{G}'(y)\) 与 \(\mathcal{G}'(x)\) 在 \(\chi\) 的一个邻域中重合。

若对每个 \(\chi \in \mathsf{X}'(\mathrm{A})\) 以及每个 \(x \in \mathrm{A}\)（\(\mathcal{G}'(x)\) 在 \(\chi\) 处为零），存在 A 中的序列 \((x_n)\) 使 \(x = \lim_{n \to \infty} x_n x\)，且每个 \(\mathcal{G}'(x_n)\) 在 \(\mathrm{V}_n\)（\(\chi\) 的一个邻域）中为零，则称 A 满足 Ditkin 条件。

注记. --- 设 A 是一个交换 Banach 代数。

\begin{enumerate}
\def\labelenumi{\arabic{enumi})}
\item
  若 \(\chi\) 使 \(\mathcal{G}'(x)\) 在 \(\chi\) 的一个邻域中为零，则 \(x\) 在 \(\chi\) 的邻域内属于 I。
\item
  若 x 在 \(\chi\) 的邻域内属于 I，且 \(y \in A\) 是任意元素，则 xy 在 \(\chi\) 的邻域内属于 I。
\item
  由特征标 \(\chi\) 组成的、使 \(x\) 在 \(\chi\) 的邻域内属于 I 的集合在 \(\mathrm{X}'(\mathrm{A})\) 中是开的。
\item
  假设 A 是正则且无根的。若 x 在每个特征标 \(\chi\) 的邻域内属于 I，其中 \(\chi \in \mathsf{X}'(\mathrm{A})\)，则 x 属于 I（I, p.~91 的推论 2 应用于函数 \(f = \mathcal{G}'(x)\)，以及 I, p.~38 的命题 8）。
\item
  假设 A 是正则的。设 I 是 A 的理想，并设 \(\chi\) 是 \(\mathsf{X}(\mathsf{A})\) 中使 \(\chi\notin\mathsf{V}(\mathsf{I})\) 的元素。那么 A 的每个元素 \(x\) 都在 \(\chi\) 的邻域内属于 I。事实上，依 I, p.~89 的定义 1，存在 \(z\in\mathsf{A}\) 使 \(\mathcal{G}'(z)\) 在 \(\chi\) 的一个邻域中等于 1，而在 V(I) 的一个邻域中等于 0。\(\mathcal{G}(z)\) 的支集是紧的，于是有 \(z\in\mathsf{I}\)（命题 4 应用于 V(I)）；因此 \(xz\in\mathsf{I}\)，且 \(\mathcal{G}'(xz)=\mathcal{G}'(x)\) 在 \(\chi\) 的一个邻域中成立。
\end{enumerate}

回顾：拓扑空间 X 的子空间 K 称为完全的，如果它是闭的且没有孤立点 (TG, I, p.~8)。

引理 1. --- 设 A 是一个无根的满足 Ditkin 条件的正则交换 Banach 代数。设 I 是 A 的一个闭理想，并设 x 是 \(\Upsilon(\mathrm{V}(\mathrm{I}))\) 的一个元素。设 K 是由 \(\chi \in \mathsf{X}'(\mathrm{A})\) 中使 x 不在 \(\chi\) 的邻域内属于 I 的特征标组成的集合。那么集合 K 是 \(\mathsf{X}'(\mathrm{A})\) 的一个完全子集。

用 G 表示 K 在 \(X'(A)\) 中的余集。集合 G 在 \(X'(A)\) 中是开的（注 3），故 K 是闭的。

用反证法，设 K 有一个孤立点 \(\chi_{0}\)。取 \(\chi_{0}\) 的一个使 \(U - \{\chi_{0}\}\subset G\) 的邻域 U。由于 x 不在 \(\chi_{0}\) 的邻域内属于 I，注 5 表明 \(\chi_{0}\in\mathrm{V(I)}\)。特别地，我们有 \(\chi_{0}(x)=0\)，因为 \(x\in\Upsilon(\mathrm{V(I)})\)。

我们将证明：存在 A 的元素 \(y\)，它在 \(\mathsf{X}'(\mathrm{A}) - \{\chi_0\}\) 的每一点的邻域内属于 I，不在 \(\chi_0\) 的邻域内属于 I，并且 \(\chi_0(y) = 0\)。

首先假设这样的元素 \(y\) 的存在性已经证明。由于 A 满足 Ditkin 条件，于是存在 A 中的序列 \((x_n)\) 使 \(x_n y\) 趋于 \(y\)，且每个 \(\mathcal{G}'(x_n)\) 在 \(\chi_0\) 的一个邻域中为零。于是对每个 \(n\)，元素 \(x_n y\) 在 \(X'(A)\) 的每一点的邻域内属于 I（注 1 与注 2），因此 \(x_{n}y \in I\)（注 4）。由于 I 是闭的，我们推出 \(y \in I\)，这与 y 不在 \(\chi_{0}\) 的邻域内属于 I 相矛盾。

剩下的就是证明 \(y\) 的存在性。若 \(\chi_0 \neq 0\)，依 I, p.~88 命题 1 的断言 (iii)，存在 \(u \in A\) 使 \(\mathcal{G}'(u)\) 在 \(\chi_0\) 的一个邻域中等于 1，而在 \(\mathrm{X}'(\mathrm{A}) - \mathrm{U}\) 的一个邻域中等于 0。设 \(y = ux\)。由于 \(x\) 对每个 \(\chi\)（这里 \(\chi \in \mathrm{U} - \{\chi_0\}\)）都在其邻域内属于 I，故对 \(y\) 亦然。此外，若 \(\chi \in \mathrm{X}'(\mathrm{A}) - \mathrm{U}\)，则 \(\mathcal{G}'(y)\) 在 \(\chi\) 的一个邻域中为零。于是（注 5）元素 \(y = ux\) 在每个 \(\chi \neq \chi_0\) 的邻域内属于 I。由于 \(\mathcal{G}'(y)\) 与 \(\mathcal{G}'(x)\) 在 \(\chi_0\) 的一个邻域中重合，\(\chi_0\) 属于 K 这一事实蕴含 \(y\) 不在 \(\chi_0\) 的邻域内属于 I。最后，我们有 \(\chi_0(y) = \chi_0(u)\chi_0(x) = 0\)。

若 \(\chi_0 = 0\)，类似地存在元素 \(v \in \mathrm{A}\) 使 \(\mathcal{G}'(v)\) 在 \(\chi_0\) 的一个邻域中为零，而在 \(\mathsf{X}'(\mathsf{A}) - \mathsf{U}\) 的一个邻域中等于 1；与前同理，我们由此推出元素 \(y = x - vx\) 在每个 \(\chi \neq \chi_0\) 的邻域内属于 I，它不在 \(\chi_0\) 的邻域内属于 I，且 \(\chi_0(y) = 0\)。

引理 2. — 设 X 是一个拓扑空间。设 F 与 D 是 X 的两个不相交的子空间，使得 F 是闭的而 D 是离散的。若 F 不含非空完全子空间，则 F ∪ D 亦然。

事实上，设 K 是 F ∪ D 的一个非空完全子空间。设 x 是 K 的一点。若 x 属于 D，则它在 D 中是孤立的，又因 F 是闭的，它在 F ∪ D 中也是孤立的。于是 x 在 K 中是孤立的，这与假设矛盾。

命题 5. — 设 A 是一个无根的满足 Ditkin 条件的正则交换 Banach 代数。设 I 是 A 的一个闭理想，使得 V(I) 的边界 F 不含非空完全集。那么 I = Y(V(I))，也就是说，I 是由使 G(x) 在 V(I) 上为零的 x ∈ A 组成的集合。特别地，若 V(I) 化为单点 χ，则 I = Ker(χ)。

我们有 \(\mathrm{I} \subset \Upsilon(\mathrm{V}(\mathrm{I}))\)。现设 \(x \in \Upsilon(\mathrm{V}(\mathrm{I}))\)。设 \(\mathrm{G}\) 是由特征标 \(\chi \in \mathsf{X}'(\mathbf{A})\) 组成的、使 \(x\) 属于 \(\mathrm{I}\)（在 \(\chi\) 的一个邻域内）的集合。它是开的，其余集 \(\mathrm{K}\) 是完全集（引理 1）。由于 \(\mathcal{G}'(x)\) 在 \(\mathrm{V}(\mathrm{I})\) 上为零，集合 \(\mathrm{G}\) 包含 \(\mathrm{V}(\mathrm{I}) \cup \{0\}\) 的内部（注 1）。依注 5，它还包含 \(\mathsf{X}(\mathbf{A}) - \mathsf{V}(\mathbf{I})\)。因此 \(\mathrm{K} = \mathsf{X}'(\mathbf{A}) - \mathsf{G}\) 含于 \(\mathrm{F}_0\)（即 \(\mathrm{V}(\mathrm{I}) \cup \{0\}\) 的边界）之中。

我们有 \(F_{0} \subset F \cup \{0\}\)。因此假设蕴含 \(F_{0}\) 不含非空完全集（引理 2）。于是由引理 1 可知完全集 K 是空集。从而 x 在每个 \(\chi \in \mathsf{X}'(\mathsf{A})\) 的邻域内属于 I，这意味着 \(x \in I\)（注 4）。这样就有 \(\Upsilon(\mathrm{V}(\mathrm{I})) \subset \mathrm{I}\)，证明完毕。

\section{星代数}

在本节中，基域是 C。

\subsection{半线性对合}

设 E 是一个复向量空间。E 上的半线性对合是从 E 到 E 的一个 R-线性映射，满足 \(u \circ u = Id_{E}\) 与 \(u(\lambda x) = \overline{\lambda}u(x)\) 对一切 \(\lambda \in C\) 及每个 \(x \in E\) 成立。这时我们用 \(E^{u}\) 表示由 \(x \in E\) 中使 \(u(x) = x\) 的元素组成的 E 的实向量子空间。

引理 1. --- 设 E 是一个复向量空间，u 是 E 上的一个半线性对合。设 \(x \in \mathrm{E}\)；令

\[
x _ {1} = \frac {1}{2} (x + u (x)), \quad x _ {2} = \frac {1}{2 i} (x - u (x)).
\]

二元组 \((x_{1}, x_{2})\) 是 \(E^{u} \times E^{u}\) 中使 \(x = x_{1} + ix_{2}\) 的唯一元素。元素 \(x_{1}\) 与 \(x_{2}\) 满足 \(x_{1} + ix_{2} = x\)，且属于 \(E^{u}\)，因为 \(u(u(x)) = x\)。反之，若 \(y_{1}\) 与 \(y_{2}\) 属于 \(E^{u}\) 且满足 \(x = y_{1} + iy_{2}\)，则得 \(u(x) = u(y_{1}) + u(iy_{2}) = y_{1} - iy_{2}\)，于是

\[
y _ {1} = \frac {1}{2} (x + u (x)) = x _ {1}, \quad i y _ {2} = \frac {1}{2} (x - u (x)) = i x _ {2}.
\]

命题 1. --- 设 \(E_{1}\) 与 \(E_{2}\) 是复向量空间，\(u_{1}\) 与 \(u_{2}\) 分别是 \(E_{1}\) 与 \(E_{2}\) 上的半线性对合。设 f 是从 \(E_{1}\) 到 \(E_{2}\) 中的一个线性映射。那么 \(f \circ u_{1} = u_{2} \circ f\) 当且仅当 \(f(\mathrm{E}_{1}^{u_{1}}) \subset \mathrm{E}_{2}^{u_{2}}\)。

若 \(f \circ u_1 = u_2 \circ f\)，我们立即得到 \(f(\mathrm{E}_1^{u_1}) \subset \mathrm{E}_2^{u_2}\)。反之，设此条件成立。设 \(x \in \mathrm{E}\)。将 x 写成 \(x = x_1 + ix_2\)，其中 \((x_1, x_2) \in \mathrm{E}_1^{u_1} \times \mathrm{E}_1^{u_1}\)（上引理）。于是有 \(f(u_1(x)) = f(x_1) - if(x_2)\) 及 \(u_2(f(x)) = u_2(f(x_1)) - iu_2(f(x_2)) = f(x_1) - if(x_2) = f(u_1(x))\)。

\subsection{对合代数}

定义 1. --- 设 A 是 C 上的一个代数。A 上的对合是从 A 到 A 的一个映射 \(x \mapsto x^{*}\)，满足：

\[
\begin{array}{c} (x ^ {*}) ^ {*} = x, \quad (x + y) ^ {*} = x ^ {*} + y ^ {*}, \quad (\lambda x) ^ {*} = \overline {{\lambda}} x ^ {*} \\ (x y) ^ {*} = y ^ {*} x ^ {*} \end{array}
\]

对一切 \(x, y \in A\) 及 \(\lambda \in C\) 成立。C 上装备了一个对合的代数称为对合代数。

A 上的对合特别地是从环 A 到反环 A° 上的一个同构。

设 A 是一个对合代数。称 \(x^{*}\) 为 x 的伴随元。A 的在对合下稳定的子集称为自伴的。若 A 有单位元 e，则有 \(e^{*} = e\)；称 (A, e) 为幺对合代数。幺对合代数的元素 u 称为酉元，如果 \(uu^{*} = u^{*}u = e\)，换言之，如果 u 可逆且其逆为 \(u^{*}\)。

元素 \(x \in A\) 称为 Hermite 元，如果 \(x = x^{*}\)；称为正规元，如果 \(xx^{*} = x^{*}x\)。这一术语推广了 A, IX, § 7, 第 3 小节中的术语。每个 Hermite 元都是正规元，每个酉元也都是正规元。A 的 Hermite 元组成的集合 \(A_{h}\) 是 A 的一个实向量子空间。若 x 与 y 是 Hermite 的且可交换，则有 \((xy)^{*} = y^{*}x^{*} = yx = xy\)，故 xy 是 Hermite 的。对每个 \(x \in A\)，A 中的元素 \(xx^{*}\) 与 \(x^{*}x\) 都是 Hermite 的。

若 A = C 装备对合 \(z \mapsto \overline{z}\)，则有 \(A_{h} = R\)。

引理 2. --- 设 A 是一个对合代数，\(x \in \mathrm{A}\)。元素

\[
x _ {1} = \frac {1}{2} (x + x ^ {*}), \qquad x _ {2} = \frac {1}{2 i} (x - x ^ {*})
\]

是 Hermite 的且满足 \(x = x_{1} + ix_{2}\)。若 \(x = y_{1} + iy_{2}\)，其中 \(y_{1}\) 与 \(y_{2}\) 是 Hermite 的，则 \(x_{1} = y_{1}\) 且 \(x_{2} = y_{2}\)。此外，元素 x 是正规元当且仅当 \(x_{1}\) 与 \(x_{2}\) 可交换。

前两个断言由 I, p.~95 的引理 1 得出。经计算有 \(xx^{*} - x^{*}x = 2i(x_{2}x_{1} - x_{1}x_{2})\)，故 x 是正规元当且仅当 \(x_{1}\) 与 \(x_{2}\) 可交换。

设 A 是一个幺对合代数。为使 \(x \in A\) 可逆，必须且只需 \(x^{*}\) 可逆，且此时有 \((x^{*})^{-1} = (x^{-1})^{*}\)。由于 \((x - \lambda e)^{*} = x^{*} - \overline{\lambda}e\) 对一切 \(\lambda \in C\) 成立，我们推出 \(\mathrm{Sp}_{\mathrm{A}}(x^{*}) = \overline{\mathrm{Sp}_{\mathrm{A}}(x)}\)。

设 A 是一个对合代数，\(\widetilde{\mathrm{A}}\) 是由 A 添加单位元而得的代数。\(\widetilde{\mathrm{A}}\) 上存在唯一的拓展 A 的对合的对合，它由 \((\lambda, x)^{*} = (\overline{\lambda}, x^{*})\) 给出，其中 \(\lambda \in \mathbf{C}\)，\(x \in \mathrm{A}\)。若 \(x \in \mathrm{A}\)，则有 \(\mathrm{Sp}_{\mathrm{A}}'(x^{*}) = \overline{\mathrm{Sp}_{\mathrm{A}}'(x)}\)。

设 A 与 B 是对合代数。从 A 到 B 的态射是从 A 到 B 的一个代数态射 \(\varphi\)，使得 \(\varphi(x^{*}) = \varphi(x)^{*}\) 对 A 中一切 \(x\) 成立。A 的恒同映射是对合代数态射。若 C 是对合代数且 \(\pi: B \to C\) 是对合代数态射，则 \(\pi \circ \varphi\) 是对合代数态射。若 \(\varphi\) 是代数同构，则 \(\varphi^{-1}\) 是对合代数态射，此时称 \(\varphi\) 为对合代数同构。

依 I, p.~95 的命题 1，若 A 与 B 是对合代数，则从 A 到 B 的代数态射 \(\varphi\) 是对合代数态射当且仅当 \(\varphi(A_h) \subset B_h\)。我们把 A 的对合子代数称为自伴子代数。A 的中心是一个对合子代数。若 \(A_1\) 是 A 的自伴双边理想，则 A 的对合通过过渡到商在代数 \(A/A_1\) 上定义一个对合，且从 A 到 \(A/A_1\) 上的典范映射是对合代数态射。

设 A 是一个对合代数。A 的根等于反代数的根 (A, VIII, p.~431, 命题 7)，因而是自伴的。

设 A 是一个对合代数。若 M ⊂ A 自伴，则其交换子 M' 是 A 的一个对合子代数。若 \(x \in A\)，则 \(\{x, x^{*}\}\) 的双交换子是包含 x 与 \(x^{*}\) 的对合子代数，且此子代数是交换的当且仅当 x 是正规元 (I, p. 5, 第 5 小节)。

注. — 设 A 是一个对合代数，B 是 A 的一个极大交换对合子代数。那么 B 是一个极大交换子代数。特别地，若 A 是幺的，则代数 B 是全的。

事实上，设 \(x \in A\) 是与 B 可交换的元素。那么 \(x^{*}\) 与 B 可交换。将 x 写成 \(x = x_{1} + ix_{2}\)，其中 \(x_{1}\) 与 \(x_{2}\) 是 Hermite 的；元素 \(x_{1}\) 与 \(x_{2}\) 与 B 可交换（引理 2）。于是 A 的由 B 与 \(x_{1}\) 生成的子代数是交换的且是对合的。从而它等于 B，故 \(x_{1} \in B\)。类似地有 \(x_{2} \in B\)，最终有 \(x \in B\)。

设 A 是一个对合代数。若 \(f\) 是 A 上的一个线性型，则 A 上的映射 \(x \mapsto \overline{f(x^*)}\) 是 A 上的线性型，记作 \(f^*\)。映射 \(f \mapsto f^*\) 是 A' 上的一个半线性对合。我们称 \(f\) 是 Hermite 的，如果 \(f = f^*\)。依 I, p.~95 的引理 1，A 上每个线性型 \(f\) 都有唯一的表示 \(f = f_1 + if_2\)，其中 \(f_1\) 与 \(f_2\) 是 Hermite 的，即 \(f_1 = \frac{1}{2}(f + f^*)\) 与 \(f_2 = \frac{1}{2i}(f - f^*)\)。

为使线性型 \(f\) 是 Hermite 的，必须且只需 \(f\) 在 \(A_h\) 上的限制取实值 (I, p.~95, 命题 1)。映射 \(f \mapsto f|A_h\) 是从 Hermite 线性型组成的实向量空间到实向量空间 \(A_h\) 的对偶向量空间上的一个同构。

特别地，我们用 \(\mathsf{X}^{\prime}(\mathsf{A})_{h}\)（相应地，\(\mathsf{X}(\mathsf{A})_{h}\)）表示 A 的 Hermite 特征标组成的集合（相应地，A 的非零 Hermite 特征标组成的集合）。于是，一个特征标是 Hermite 的，如果它在 \(A_{h}\) 上的限制取实值。

若 A 是交换的且 \(\chi\) 是 A 的特征标，则 \(\chi^{*}\) 是 A 的特征标，且映射 \(\chi \mapsto \chi^{*}\) 是从 \(\mathsf{X}'(\mathsf{A})\) 到 \(\mathsf{X}'(\mathsf{A})\) 上的一个同胚。

引理 3. --- 设 A 是一个交换对合代数。A 到 \(\mathcal{C}_{0}(\mathsf{X}(\mathsf{A}))\) 的 Gelfand 变换是对合代数态射当且仅当 A 的每个特征标都是 Hermite 的。

事实上，说 \(G_{A}\) 是对合代数态射，等于说对一切 \(x \in A\) 与 \(\chi \in \mathsf{X}(A)\) 有

\[
\chi (x ^ {*}) = \mathcal {G} _ {\mathrm{A}} (x ^ {*}) (\chi) = \overline {{\mathcal {G} _ {\mathrm{A}} (x) (\chi)}} = \overline {{\chi (x)}},
\]

也就是说，每个 \(\chi\) 都是 Hermite 的。

例. --- 1) 设 A 是集合 X 上复值函数组成的代数。映射 \(f \mapsto \overline{f}\) 是 A 中的一个对合。X 上有界函数组成的子代数是 A 的对合子代数。若 X 是局部紧拓扑空间，则子代数 \(\mathcal{C}(X), \mathcal{C}_b(X), \mathcal{C}_0(X)\) 与 \(\mathcal{K}(X)\) 都是 A 的对合子代数。

\begin{enumerate}
\def\labelenumi{\arabic{enumi})}
\setcounter{enumi}{1}
\item
  设 X 是局部紧拓扑空间，\(\mu\) 是 X 上的一个正测度。映射 \(f \mapsto \overline{f}\) 是代数 \(\mathcal{L}^{\infty}(X, \mu)\) 上的一个对合；它通过过渡到商在幺代数 \(\mathrm{L}^{\infty}(X, \mu)\) 上诱导出一个对合。
\item
  设 E 是一个复 Hilbert 空间。在 Banach 代数 \(\mathcal{L}(\mathrm{E})\) 上，映射 \(x \mapsto x^{*}\) (EVT, V, p.~37, 命题 1) 是一个对合。
\item
  设 G 是一个局部紧群。设 \(\mathcal{M}^{1}(G)\) 是 G 上有界复测度组成的 Banach 代数 (I, p.~19, 例 6)。
\end{enumerate}

从 G 到 G 的映射 \(x \mapsto x^{-1}\) 把每个测度 \(\mu \in \mathcal{M}^{1}(G)\) 映为测度 \(\check{\mu} \in \mathcal{M}^{1}(G)\) (INT, VII, p.~12, 公式 (13))。我们用 \(\mu^{*}\) 表示 \(\check{\mu}\) 的复共轭测度。映射 \(\mu \mapsto \check{\mu}\) 是从 Banach 代数 \(\mathcal{M}^{1}(G)\) 到 Banach 代数 \(\mathcal{M}^{1}(G^{\circ})\) 上的一个等距同构 (INT, VIII, §3, 第 1 小节, 命题 7 的推论)；于是 \(\mu \mapsto \mu^{*}\) 是 Banach 代数 \(\mathcal{M}^{1}(G)\) 的一个等距对合。

关于一个 Haar 测度具有密度的有界测度组成的集合 A 是 \(\mathcal{M}^{1}(\mathrm{G})\) 的一个在对合下稳定的闭子代数（参见 INT, VIII, §4, 第 5 小节）；它与 Haar 测度的选取无关。

设 \(\nu\) 是 G 上的一个左 Haar 测度，用 \(\Delta\) 表示 G 的模函数。我们给 \(\mathrm{L}^1 (\mathrm{G},\nu)\) 装备乘积 \((f,g)\mapsto f*\nu g\) 与对合 \(f\mapsto f^{*} = \widetilde{f}\cdot \Delta^{-1}\)，其中 \(\widetilde{f} (x) = \overline{f(x^{-1})}\) 对一切 \(x\in \mathbf{G}\) 成立。于是映射 \(f\mapsto f\cdot \nu\) 是从对合代数 \(\mathrm{L}^1 (\mathrm{G},\nu)\) 到 A 上的一个同构。这个同构是等距的。特别地，\(\mathrm{L}^1 (\mathrm{G},\nu)\) 恒等同于 \(\mathcal{M}^1 (\mathrm{G})\) 的一个对合子代数。

\begin{enumerate}
\def\labelenumi{\arabic{enumi})}
\setcounter{enumi}{4}
\tightlist
\item
  设 U 是 C 的一个在复共轭下稳定的开子集。考虑 U 上复值全纯函数组成的代数 \(\mathcal{O}(\mathrm{U})\)。对每个函数 \(f \in \mathcal{O}(\mathrm{U})\)，映射 \(f^{*}: z \mapsto \overline{f(\overline{z})}\) 是 U 上的一个全纯函数。映射 \(f \mapsto f^{*}\) 是 \(\mathcal{O}(\mathrm{U})\) 上的一个对合。
\end{enumerate}

类似地，设 S 是 C 的一个在复共轭下稳定的紧子集。考虑在 S 的邻域中的复值全纯函数芽组成的代数 \(\mathcal{O}(S)\)。映射 \(f \mapsto f^{*}\) 是 \(\mathcal{O}(S)\) 上的一个对合。

在 I, p.~85 第 13 小节中定义的子代数 \(\mathcal{O}_{\mathbf{R}}(\mathrm{U})\)（相应地，子代数 \(\mathcal{O}_{\mathbf{R}}(\mathrm{S})\)）是 \(\mathcal{O}(\mathrm{U})\)（相应地，\(\mathcal{O}(\mathrm{S})\)）的 Hermite 元组成的集合。

\subsection{赋范对合代数}

定义 2. --- 赋范对合代数是装备了一个对合 \(x \mapsto x^{*}\) 的赋范代数 A，满足 \(\|x^{*}\| = \|x\|\) 对一切 \(x\) 成立。若 A 是 Banach 代数，则称 A 为对合 Banach 代数。

例. --- 1) 设 X 是局部紧拓扑空间。X 上连续有界复函数组成的 Banach 代数 \(\mathcal{C}_{b}(\mathrm{X})\)，装备范数 \(\|f\| = \sup_{x\in\mathrm{X}}|f(x)|\) 与对合 \(f\mapsto\overline{f}\)，是一个对合 Banach 代数。在无穷远处趋于 0 的连续函数组成的子代数 \(\mathcal{C}_{0}(\mathrm{X})\) 是 \(\mathcal{C}_{b}(\mathrm{X})\) 的一个闭对合子代数。

\begin{enumerate}
\def\labelenumi{\arabic{enumi})}
\setcounter{enumi}{1}
\item
  设 X 是局部紧拓扑空间，\(\mu\) 是 X 上的一个正测度。对合代数 \(\mathrm{L}^{\infty}(\mathrm{X},\mu)\) (I, p.~99, 例 2) 是对合 Banach 代数，因为 \(|f| = |\overline{f}|\) 对每个元素 \(f \in \mathrm{L}^{\infty}(\mathrm{X},\mu)\) 成立。
\item
  复 Hilbert 空间 E 的连续自同态组成的对合代数 \(\mathcal{L}(\mathrm{E})\) (I, p.~99, 例 3)，装备范数
\end{enumerate}

\[
\| u\| = \sup_{\substack{x\in \operatorname{E}\\ \| x\| \leqslant 1}}\| u(x)\|
\]

（EVT, III, p.~14）后，是一个对合 Banach 代数 (EVT, V, p.~37, 命题 1)。

\begin{enumerate}
\def\labelenumi{\arabic{enumi})}
\setcounter{enumi}{3}
\item
  局部紧群上有界测度组成的对合代数 \(\mathcal{M}^{1}(G)\) (I, p.~99, 例 4)，装备通常的范数 (I, p.~19, 例 6)，是对合 Banach 代数。设 \(\nu\) 是 G 上的一个左 Haar 测度。对合 Banach 代数 \(\mathrm{L}^{1}(\mathrm{G},\nu)\) 恒等同于 \(\mathcal{M}^{1}(\mathrm{G})\) 的一个闭子代数。
\item
  设 \((\mathrm{A}_{i})\) 是一族对合赋范代数。设 A 是 \(A_{i}\) 的乘积赋范代数 (I, p.~15, 第 1 小节)。代数 A 装备对合 \((x_{i})^{*} = (x_{i}^{*})\) 后是一个对合赋范代数。若每个代数 \(A_{i}\) 都是对合 Banach 代数，则 A 是对合 Banach 代数。我们称 A 为 \(A_{i}\) 的乘积对合赋范代数（相应地，乘积对合 Banach 代数）。
\item
  设 A 是一个赋范对合代数，\(\widetilde{\mathrm{A}}\) 是由 A 添加单位元而得的赋范代数。代数 \(\widetilde{\mathrm{A}}\) 装备第 2 小节中定义的对合后是一个赋范对合代数。若 A 是对合 Banach 代数，则 \(\widetilde{\mathrm{A}}\) 也是对合 Banach 代数。
\end{enumerate}

若 A 是赋范对合代数，对合子代数的闭包是对合子代数。若 M ⊂ A，包含 M 的最小闭对合子代数称为由 M 生成的闭对合子代数；它是由 M ∪ M* 生成的子代数的闭包。若 M 由单个正规元组成，则由 M 生成的闭对合子代数是交换的，且其所有元素都是正规的。

类似地，若 A 是幺对合赋范代数且 M 是 A 的一个子集，包含 M 的最小闭幺对合子代数称为由 M 生成的闭幺对合子代数；它是由 M ∪ M* 生成的幺子代数的闭包。若 M 由单个正规元组成，则由 M 生成的闭幺对合子代数是交换的，且其所有元素都是正规的。

对合赋范代数对闭自伴双边理想的商、完备化以及对合赋范代数的反代数，都自然是赋范对合代数。

若 A 是赋范对合代数，A 的 Hermite 元组成的集合 \(A_{h}\) 是一个实赋范向量空间。

引理 4. --- 设 A 是一个赋范对合代数。对 A 上每个连续线性型 f，我们有 \(\|f^{*}\|=\|f\|\)。此外若 f 是 Hermite 的，则 \(\|f\|=\|f|A_{h}\|\)。

第一个断言由定义得出。对第二个断言，设 \(g\) 是 \(f\) 在 \(\mathrm{A}_h\) 上的限制。我们有 \(\| f \| \geqslant \| g \|\)。我们来证明反向不等式。对每个 \(\varepsilon > 0\)，存在 \(x \in \mathrm{A}\) 使 \(\| x \| \leqslant 1\) 且 \(|f(x)| \geqslant \|f\| - \varepsilon\)。用模为 1 的复数乘 x，可设 \(f(x) \geqslant 0\)。于是元素 \(\frac{1}{2}(x + x^{*})\) 属于 \(A_{h}\) 且范数 \(\leqslant 1\)。我们有

\[
\left| g \left(\frac {1}{2} (x + x ^ {*})\right) \right| = \frac {1}{2} | f (x) + f (x ^ {*}) | = f (x) \geqslant \| f \| - \varepsilon
\]

于是 \(\|g\|\geqslant\|f\|-\varepsilon\)。由此我们推出 \(\|g\|\geqslant\|f\|\)。

在下文中，我们将把 A 上连续的 Hermite 线性型与 \(A_{h}\) 上连续的实线性型恒同。

引理 5. --- 设 A 是一个对合 Banach 代数。

\begin{enumerate}
\def\labelenumi{\alph{enumi})}
\item
  对每个 \(x \in A\)，有 \(\exp(x)^{*} = \exp(x^{*})\)；
\item
  设 \(x \in \mathrm{A}_h\) 是一个 Hermite 元。那么 \(\exp(ix)\) 是酉元。
\end{enumerate}

事实上，由于 A 上的对合是连续的，我们有

\[
\exp (x) ^ {*} = \left(\sum_ {n = 0} ^ {\infty} \frac {x ^ {n}}{n !}\right) ^ {*} = \sum_ {n = 0} ^ {\infty} \frac {(x ^ {*}) ^ {n}}{n !} = \exp (x ^ {*})
\]

对一切 \(x \in A\) 成立 (I, p.~78, 公式 (9))。若 \(x \in A_{h}\)，则由此得

\[
\exp (i x) ^ {*} = \exp ((i x) ^ {*}) = \exp (- i x) = \exp (i x) ^ {- 1}
\]

(I, p.~78, 公式 (11))。

\subsection{星代数}

定义 3. --- 星代数（C*-代数）是这样的对合 Banach 代数 A：\(\|x\|^{2}=\|x^{*}x\|\) 对一切 \(x\in A\) 成立。

若 A 与 B 是星代数，从 A 到 B 的态射，或称星代数态射，是指从 A 到 B 的对合代数态射。从 A 到 B 的同构是指从 A 到 B 的对合代数同构。

某些作者称之为“C*-代数”。

设 A 是一个星代数。A 的星子代数是指 A 的一个闭对合子代数。

例. --- 1) 复 Hilbert 空间的连续自同态组成的对合 Banach 代数 (I, p.~100, 例 3) 是一个星代数 (EVT, V, p.~39, 命题 2)。

\begin{enumerate}
\def\labelenumi{\arabic{enumi})}
\setcounter{enumi}{1}
\tightlist
\item
  设 X 是局部紧拓扑空间。X 上连续有界复值函数组成的对合 Banach 代数 \(\mathcal{C}_{b}(X)\) (I, p.~100, 例 1) 是一个星代数。事实上，对每个函数 \(f \in \mathcal{C}_b(\mathrm{X})\)，有 \(f^*f = |f|^2\)，从而 \(\| f^*f \| = |||f|^2 \| = \| f\|^2\)。
\end{enumerate}

设 \(\mathrm{A} = \mathcal{C}_0(\mathrm{X})\) 是在无穷远处趋于 0 的函数组成的对合 Banach 子代数。它是 \(\mathcal{C}_b(\mathrm{X})\) 的一个星子代数。对每个函数 \(f \in \mathrm{A}\)，有 \(\| f \| = \varrho(f)\)，因为 \(\mathrm{Sp}_{\mathrm{A}}'(f) = f(\mathrm{X}) \cup \{0\}\)。

设 X 与 Y 是局部紧拓扑空间。对每个从 X 到 Y 中的真偏映射 \(\varphi\) (I, p.~33, 定义 1)，代数态射 \(\varphi^{*}\)（从 \(\mathcal{C}_{0}(\mathrm{Y})\) 到 \(\mathcal{C}_{0}(\mathrm{X})\)，I, p.~34, 命题 3）是对合代数态射。反之，每个星代数态射 \(\pi\colon\mathcal{C}_{0}(\mathrm{Y})\to\mathcal{C}_{0}(\mathrm{X})\) 都具有这种形式（同前）。

\begin{enumerate}
\def\labelenumi{\arabic{enumi})}
\setcounter{enumi}{2}
\item
  设 X 是紧拓扑空间，\(x_{0} \in X\) 是 X 的一个固定元素。\(\mathcal{C}'(X)\) 是 \(\mathcal{C}(X)\) 的一个由连续函数 \(f: X \to C\)（其中 \(f(x_{0}) = 0\)）组成的子代数，它是一个星代数。
\item
  设 X 是分离拓扑空间，\(\mu\) 是 X 上的一个正测度。对合 Banach 代数 \(\mathrm{L}^{\infty}(\mathrm{X},\mu)\) 是一个交换的幺星代数。
\item
  设 \((\mathbf{A}_{i})\) 是一族星代数。\(A_{i}\) 的乘积对合 Banach 代数 A (I, p.~101, 例 5) 是一个星代数，称为 \(A_{i}\) 的乘积星代数。
\item
  设 A 是一个星代数。若 B 是 A 的一个闭对合子代数，则 B 是一个星代数。（V，将出版）中将看到，每个星代数都同构于一个 Hilbert 空间的自同态组成的星代数（例 1）的一个闭对合子代数。
\item
  设 A 是一个星代数。若 M ⊂ A 是任意子集，则 A 的由 M 生成的闭对合子代数是一个星代数，称为 A 的由 M 生成的星子代数。此外若 A 是幺的，则由 M 生成的闭幺对合子代数是一个幺星代数，称为 A 的由 M 生成的幺星子代数。
\item
  一般而言，对合 Banach 代数 \(\mathcal{M}^{1}(G)\) (I, p.~100, 例 4) 不是星代数 (I, p.~181, 习题 8)。
\end{enumerate}

引理 6. --- 设 A 是一个 Banach 代数，装备了一个满足

\[
\| x \| ^ {2} \leqslant \| x ^ {*} x \|\tag{1}
\]

（对一切 \(x \in \mathbf{A}\)）的对合。那么 A 是一个星代数。

设 \(x \in \mathbf{A}\)。于是有 \(\|x\|^2 \leqslant \|x^*\| \cdot \|x\|\)，从而 \(\|x\| \leqslant \|x^*\|\)。交换 \(x\) 与 \(x^*\) 的角色，我们得到 \(\|x\| = \|x^*\|\)。于是 \(\|x^*x\| \leqslant \|x^*\|\|x\| \leqslant \|x\|^2\)，而假设蕴含等式 \(\|x\|^2 = \|x^*x\|\)。

引理 7. --- 设 A 是一个星代数。

\begin{enumerate}
\def\labelenumi{\alph{enumi})}
\tightlist
\item
  A 的正则表示 \(\gamma\) (I, p.~16, 定义 1) 是等距的，也就是说，
\end{enumerate}

\[
\| x \| = \sup _ {\| y \| \leqslant 1} \| x y \|,
\]

对一切 \(x\in \mathbf{A}\) 成立。

\begin{enumerate}
\def\labelenumi{\alph{enumi})}
\setcounter{enumi}{1}
\tightlist
\item
  对每个 \(x \in A_{h}\)，有
\end{enumerate}

\[
\varrho (x) = \| x \|.\tag{2}
\]

设 \(x \in \mathrm{A}\)。我们有 \(\sup_{\| y \| \leqslant 1} \| xy \| \leqslant \| x \|\)。为证 \(\| x \| \leqslant \sup_{\| y \| \leqslant 1} \| xy \|\)，可设 \(\| x \| = 1\)。于是元素 \(y = x^*\) 满足 \(\| y \| = \| x^* \| = 1\)，且 \(\| xy \| = \| x \|^2 = 1\)，由此得断言 \(a\))。

设 \(x\) 是 Hermite 的。由此得 \(\| x^2 \| = \| x^* x \| = \| x \|^2\)，从而用归纳法得 \(\| x^{2^n} \|^{2^{-n}} = \| x \|\) 对每个整数 \(n \geqslant 1\) 成立，再依 I, p.~20 的命题 1 即得断言 \(b)\)。

注记. --- 1) 设 A 是一个幺星代数。那么我们有

\[
\left\| 1 \right\| ^ {2} = \left\| 1 ^ {*} 1 \right\| = \left\| 1 \right\|,
\]

因此范数 \(\|1\|\) 或为零或等于 1。若 A ≠ \{0\}，我们推出 \(\|1\| = 1\)。从而对 A 的每个酉元 u，有 \(\|u\| = \|u^{*}u\|^{1/2} = 1\)。

\begin{enumerate}
\def\labelenumi{\arabic{enumi})}
\setcounter{enumi}{1}
\tightlist
\item
  设 A 是一个赋范对合代数。若 \(\|x\|^{2}=\|x^{*}x\|\) 对一切 \(x\in A\) 成立，则 A 的完备化 \(\widehat{A}\) 是一个星代数。
\end{enumerate}

命题 2. --- 设 A 是一个对合 Banach 代数，B 是一个星代数，π 是从 A 到 B 的一个对合代数态射。那么 \(\|\pi(x)\| \leqslant \|x\|\) 对一切 \(x \in A\) 成立，特别地，π 是连续的。

对每个 \(x \in \mathrm{A}\)，有 \(\mathrm{Sp}_{\mathrm{B}}^{\prime}(\pi(x)) \subset \mathrm{Sp}_{\mathrm{A}}^{\prime}(x)\)，于是 \(\varrho(\pi(x)) \leqslant \varrho(x) \leqslant \|x\|\)。由于 \(\pi(x^{*}x) \in \mathrm{B}_{h}\)，有 \(\|\pi(x^{*}x)\| = \varrho(\pi(x^{*}x))\)（公式 (2)），从而

\[
\| \pi (x) \| ^ {2} = \| \pi (x ^ {*} x) \| = \varrho (\pi (x ^ {*} x)) \leqslant \| x ^ {*} x \| = \| x \| ^ {2}.
\]

命题 3. --- 设 A 是一个星代数，\(\widetilde{\mathrm{A}}\) 是由 A 添加单位元而得的对合代数。\(\widetilde{\mathrm{A}}\) 上存在唯一的范数，它拓展 A 的范数并使 \(\widetilde{\mathrm{A}}\) 成为一个星代数。

这样的范数的唯一性由命题 2 得出。现在证明其存在性。用 \(\widetilde{e}\) 表示 \(\widetilde{\mathrm{A}}\) 的单位元。

首先，设 A 有单位元 e。对合赋范代数 A 与 \(\mathbf{C}(\widetilde{e}-e)\) 的乘积恒等同于 \(\widetilde{A}\)，且是一个星代数（例 5）。\(\widetilde{A}\) 上的范数拓展 A 的范数，由此得断言。

现设 A 没有单位元。对每个 \(x \in \widetilde{A}\)，用 \(\gamma_{x}\) 表示从 A 到 A 的乘法算子 \(y \mapsto xy\)，并令 \(\|x\|_{\widetilde{A}} = \| \gamma_{x}\|\)。映射 \(x \mapsto \|x\|_{\widetilde{A}}\) 是 \(\widetilde{A}\) 上的一个半范数。对一切 x 与 \(x'\)（属于 \(\widetilde{A}\)），有 \(\|xx'\|_{\widetilde{A}} \leqslant \|x\|_{\widetilde{A}} \|x'\|_{\widetilde{A}}\)。此外，依引理 7，\(\|x\|_{\widetilde{A}} = \|x\|\) 对一切 \(x \in \widetilde{A}\) 成立。

我们来证明映射 \(x \mapsto \|x\|_{\widetilde{\mathrm{A}}}\) 是 \(\widetilde{\mathrm{A}}\) 上的一个范数。设 \(\lambda \in \mathbf{C}\) 与 \(x \in \mathrm{A}\) 使 \(\| \lambda \widetilde{e} + x\|_{\widetilde{\mathrm{A}}} = 0\)。若 \(\lambda \neq 0\)，则条件 \((\lambda \widetilde{e} + x)y = 0\)（对一切 \(y \in \mathrm{A}\)）蕴含 \(-\lambda^{-1}x\) 是 A 中的一个左单位元。同样，元素 \(-\overline{\lambda}^{-1}x^{*}\) 是一个右单位元。于是代数 A 就会有一个单位元，这与假设相反。故 \(\lambda = 0\)。但这时 \(0 = \|x\|_{\widetilde{\mathrm{A}}} = \|x\|\)，从而 \(x = 0\)。

由于 A 是完备的且在 \(\widetilde{\mathrm{A}}\) 中的余维数为 1，故空间 \(\widetilde{\mathrm{A}}\)（装备范数 \(x \mapsto \|x\|_{\widetilde{\mathrm{A}}}\)）是完备的。因此，只需证明 \(\|x\|_{\widetilde{\mathrm{A}}}^{2} \leqslant \|x^{*}x\|_{\widetilde{\mathrm{A}}}\) 对一切 \(x \in \widetilde{\mathrm{A}}\) 成立（引理 6）即可完成证明。可设 \(\|x\|_{\widetilde{\mathrm{A}}} = 1\)。于是对每个实数 \(r < 1\)，存在 \(y \in \mathrm{A}\) 使 \(\|y\| = \|y^{*}\| \leqslant 1\) 且 \(\|xy\|^2 \geqslant r\)。由于 \(xy \in \mathrm{A}\)，有

\[
\| x ^ {*} x \| _ {\widetilde {A}} \geqslant \| x ^ {*} x y \| \geqslant \| y ^ {*} (x ^ {*} x) y \| = \| (x y) ^ {*} (x y) \| = \| x y \| ^ {2} \geqslant r.
\]

由此推出 \(\| x^{*}x\|_{\widetilde{\mathrm{A}}}\geqslant 1\)，从而对合代数 \(\widetilde{\mathrm{A}}\)（装备范数 \(x\mapsto \| x\|_{\widetilde{\mathrm{A}}}\)）是一个星代数。

定义 4. --- 称装备命题 3 中范数的 \(\widetilde{\mathrm{A}}\) 为由 A 添加单位元而得的星代数。

当 \(A \neq \{0\}\) 时，赋范对合代数 \(\widetilde{A}\) 上的星代数范数并不是 I, p.~101 例 6 中考虑的那个（参见 I, p.~181, 习题 10）。

命题 4. --- 设 A 是一个星代数。

\begin{enumerate}
\def\labelenumi{\alph{enumi})}
\item
  若 A 有单位元且 u 是 A 的一个酉元，则 \(\mathrm{Sp}(u) \subset \mathbf{U}\)；
\item
  若 h 是 A 的一个 Hermite 元，则 \(\mathrm{Sp}'(h) \subset \mathbf{R}\)。
\end{enumerate}

可设 A 非零。我们来证明断言 \(a\))。设 \(u\) 是 A 的一个酉元。我们有 \(\|u\|=\|u^{-1}\|=1\)（注 1），于是 \(\mathrm{Sp}(u)\subset\mathbf{U}\) (I, p.~26, 推论 3)。为证 \(b\))，可设 A 是幺的（命题 3）。设 \(h\) 是 A 的一个 Hermite 元。那么 \(\exp(ih)\) 是酉元 (I, p.~102, 引理 5)。于是，依 I, p.~67 的推论 2 及 \(a\))，有 \(\exp(i\mathrm{Sp}(h))=\mathrm{Sp}(\exp(ih))\subset\mathbf{U}\)，这就意味着 \(\mathrm{Sp}(h)\subset\mathbf{R}\)。

命题 5. — 设 A 是一个幺星代数，B 是 A 的包含 A 的单位元的一个星子代数。那么 B 是 A 的一个全子代数。特别地，\(\mathrm{Sp}_{\mathrm{B}}(x) = \mathrm{Sp}_{\mathrm{A}}(x)\) 对 B 中一切 \(x\) 成立。

设 x 是 B 的一个 Hermite 元。由于 \(\mathrm{Sp}_{\mathrm{B}}(x) \subset \mathbf{R}\)（命题 4），I, p.~28 的命题 6 表明 \(\mathrm{Sp}_{\mathrm{B}}(x) = \mathrm{Sp}_{\mathrm{A}}(x)\)。特别地，x 在 B 中可逆当且仅当它在 A 中可逆。

现设 x 是 B 的任一在 A 中可逆的元素。那么 \(x^{*}\) 在 A 中可逆，且 \(xx^{*}\) 在 A 中可逆。由于 \(xx^{*}\) 是 Hermite 的，前面已证 \(xx^{*}\) 在 B 中可逆。这蕴含 x 在 B 中右可逆。类似地，可以验证 x 在 B 中左可逆，从而 x 在 B 中可逆。于是，B 是 A 的一个全子代数。

推论. --- 设 A 是一个星代数，B 是 A 的一个星子代数。那么对 B 中一切 x 有 \(\mathrm{Sp}_{\mathrm{B}}^{\prime}(x)=\mathrm{Sp}_{\mathrm{A}}^{\prime}(x)\)。

通过添加单位元（命题 3），这由命题 5 得出。

命题 6（Fuglede–Putnam 定理）

设 A 是一个幺星代数。设 a 与 b 是 A 的正规元。若 \(c \in A\) 满足 ac = cb，则 \(a^{*}c = cb^{*}\)。

假设蕴含 \((wa)^{k}c = c(wb)^{k}\) 对每个整数 \(k \geqslant 0\) 及每个 \(w \in C\) 成立，从而 \(e^{wa}c = ce^{wb}\) (I, p.~78, 公式 (9))。考虑由 \(z \mapsto e^{-za^{*}}ce^{zb^{*}}\) 定义的从 C 到 A 的函数 f。它是 C 上的一个全纯函数，其导数满足

\[
f ^ {\prime} (z) = - a ^ {*} e ^ {- z a ^ {*}} c e ^ {z b ^ {*}} + e ^ {- z a ^ {*}} c b ^ {*} e ^ {z b ^ {*}}
\]

对一切 \(z \in \mathbf{C}\) 成立。由于 \(c = e^{-\overline{z}a} ce^{\overline{z}b}\)，可以写出

\[
f (z) = e ^ {- z a ^ {*}} e ^ {\overline {{z}} a} c e ^ {- \overline {{z}} b} e ^ {z b ^ {*}}.
\]

由于 \(a\) 与 \(b\) 是正规的，A 中元素 \(e^{-za^*}e^{\overline{z} a} = e^{-za^* + \overline{z} a}\) 与 \(e^{-\overline{z} b}e^{zb^*} = e^{-\overline{z} b + zb^*}\) 对一切 \(z\in \mathbf{C}\) 都是酉的 (I, p.~102, 引理 5)，从而范数为 1。于是 \(\| f(z)\| \leqslant \| c\|\) 对一切 \(z\in \mathbf{C}\) 成立。因此函数 \(f\) 是常数 (VAR, R1, 3.3.6, p.~29)，即 \(f(z) = f(0) = c\) 对一切 \(z\in \mathbf{C}\) 成立。但这时 \(-a^{*}c + cb^{*} = f'(0) = 0\)。

推论. --- 设 A 是一个星代数，a 是 A 的一个正规元。\(\{a, a^{*}\}\) 的交换子（相应地，双交换子）与 a 的交换子（相应地，双交换子）重合。

只需在命题中取 b = a。

\subsection{交换星代数}

引理 8. --- 设 X 与 Y 是度量空间，其中 X 是完备的。设 f 是从 X 到 Y 中的一个映射，使得

\[
d (f (x), f (y)) \geqslant d (x, y)
\]

对 X 中一切 x 与 y 成立，且使得 f 的图像在 X × Y 中是闭的。那么 f 是一个闭映射。

设 F 是 X 的一个闭子集，\((y_{n})_{n\in\mathbf{N}}\) 是 \(f(\mathrm{F})\) 中收敛于 \(y\in Y\) 的一个序列。对每个 \(n\in N\)，取 \(x_{n}\in F\) 使 \(f(x_{n})=y_{n}\)。假设蕴含序列 \((x_{n})_{n\in\mathbf{N}}\) 是 Cauchy 序列；设 \(x\in X\) 是它的极限；由于 F 是闭的，它是 F 的元素。此外，我们有 \((x_{n},f(x_{n}))\to(x,y)\) 在 \(X\times Y\) 中成立。由于 f 的图像是闭的，由此得 \(y=f(x)\) 属于 \(f(\mathrm{F})\)。

引理 9. --- 设 A 是一个交换星代数。A 的每个特征标都是 Hermite 的，且 Gelfand 变换是从 A 到 \(\mathcal{C}_{0}(\mathsf{X}(\mathsf{A}))\) 中的星代数态射。

只需证明第一个断言（I, p.~38 的命题 7 与 I, p.~98 的引理 3）。设 \(\chi\) 是 A 的一个特征标。对 A 的每个 Hermite 元 \(y\)，我们有 \(\chi(y) = \mathcal{G}_{\mathrm{A}}(y)(\chi) \in \mathrm{Sp}'(y) \subset \mathbf{R}\) (I, p.~106, 命题 4)。从而特征标 \(\chi\) 是 Hermite 的 (I, p.~95, 命题 1)。

定理 1. --- 设 A 是一个交换星代数。Gelfand 变换是从星代数 A 到星代数 \(\mathcal{C}_{0}(\mathsf{X}(\mathsf{A}))\)（由 \(\mathsf{X}(\mathsf{A})\) 上在无穷远处为 0 的连续函数组成）上的一个等距同构。

Gelfand 变换是从 A 到 \(\mathcal{C}_{0}(\mathsf{X}(\mathsf{A}))\) 中的一个对合代数态射（引理 9）。设 B 是它的像。它是 \(\mathcal{C}_{0}(\mathsf{X}(\mathsf{A}))\) 的一个对合子代数。按定义，B 的元素分离 \(\mathsf{X}(\mathsf{A})\) 的点。设 \(\chi \in \mathsf{X}(\mathsf{A})\)。存在 \(x \in \mathsf{A}\) 使 \(\chi(x) \neq 0\)，于是 \(f = \mathcal{G}(x)\) 是 B 中使 \(f(\chi) \neq 0\) 的一个元素。依 TG, X, p.~40, 推论 2，子代数 B 因此在 \(\mathcal{C}_{0}(\mathsf{X}(\mathsf{A}))\) 中稠密。

对 A 的每个 Hermite 元 \(y\)，我们有 \(\| \mathcal{G}(y)\| = \varrho (y) = \| y\|\)（I, p.~38 的命题 7 与 I, p.~104 的公式 (2)），从而对每个 \(x\in \mathrm{A}\) 有等式

\[
\| x \| ^ {2} = \| x ^ {*} x \| = \| \mathcal {G} (x ^ {*} x) \| = \| \overline {{\mathcal {G} (x)}} \cdot \mathcal {G} (x) \| = \| \mathcal {G} (x) \| ^ {2}.
\]

于是 G 是等距的，因此它的像 B 是闭的（引理 8）。我们得出结论 \(\mathrm{B} = \mathcal{C}_{0}(\mathrm{X}(\mathrm{A}))\)。

推论 1. --- 设 A 是一个星代数，x 是 A 的一个正规元。那么 \(\|x\| = \varrho(x)\)。

由于 A 的由 x 与 \(x^{*}\) 生成的星子代数是交换的，可设 A 是交换的。在这种情形，推论由定理 1、I, p.~102 的例 2 以及 I, p.~24 的定理 1 得出。

推论 2. --- 设 A 是一个交换星代数。

\begin{enumerate}
\def\labelenumi{\alph{enumi})}
\item
  存在局部紧拓扑空间 X，使 A 同构于星代数 \(\mathcal{C}_{0}(\mathrm{X})\)；
\item
  设 B 是一个交换星代数。映射 \(\pi \mapsto \mathsf{X}'(\pi)\) 是从 A 到 B 的星代数态射组成的集合到从 \(\mathsf{X}(\mathsf{B})\) 到 \(\mathsf{X}(\mathsf{A})\) 中的真偏映射组成的集合上的一个双射 (I, p.~33, 定义 1)。
\end{enumerate}

定理 1 建立了第一个断言，第二个断言由 I, p.~34 的命题 3 与 I, p.~102 的例 2 得出。

推论 3. --- 设 A 是一个幺交换星代数。

\begin{enumerate}
\def\labelenumi{\alph{enumi})}
\item
  存在紧拓扑空间 X，使 A 同构于星代数 \(\mathcal{C}(\mathrm{X})\)；
\item
  设 B 是一个交换幺星代数。映射 \(\pi \mapsto \mathsf{X}(\pi)\) 是从 A 到 B 的幺星代数态射组成的集合到从 \(\mathsf{X}(\mathsf{B})\) 到 \(\mathsf{X}(\mathsf{A})\) 中的连续映射组成的集合上的一个双射。
\end{enumerate}

这由前文及 I, p.~35 的命题 4 得出。

注记. --- *设 G 为以局部紧空间为对象、以真偏映射 (I, p.~33, 定义 1) 为态射的范畴，并设 S 为交换星代数的范畴，其态射是对合代数态射。考虑从 S 到反范畴 G° 的函子，它把交换星代数 A 对应到 A 的非零特征标组成的局部紧空间 X(A)，把交换星代数的态射 φ: A → B 对应到连续映射 X'(φ)。定理 1 与推论 2 表明这个函子是一个范畴等价，且该函子的一个拟逆是这样的函子：它把局部紧拓扑空间 X 对应到交换星代数 C₀(X)。

类似地，推论 3 表明紧空间范畴的反范畴等价于幺交换星代数的范畴。*

\subsection{幺星代数中的函数演算}

在本小节中，A 是一个幺星代数，x 是 A 的一个正规元。

设 B 是 A 的由 x 生成的幺星子代数；它是交换的，且包含在 \(\{x, x^{*}\}\) 的双交换子中，从而包含在 x 的双交换子中 (I, p.~106 命题 6 的推论)。Gelfand 变换 \(\mathcal{G}_{\mathrm{B}}: \mathrm{B} \to \mathcal{C}(\mathrm{X}(\mathrm{B}))\) 是星代数的一个同构 (I, p.~108, 定理 1)。我们有 \(\mathrm{Sp}_{\mathrm{B}}(x) = \mathrm{Sp}_{\mathrm{A}}(x)\) (I, p.~106, 命题 5)。

引理 10. --- 映射 \(\mathrm{ev}_x\colon \chi \mapsto \chi (x)\) 诱导出从 \(\mathsf{X}(\mathbf{B})\) 到 \(\mathrm{Sp}_{\mathrm{A}}(x)\) 上的一个同胚。

\(x \mapsto \chi(x)\) 是从 \(\mathsf{X}(\mathsf{B})\) 到 \(\mathbf{C}\) 中的一个连续映射，依 I, p.~37 的命题 6，它的像等于 \(\mathrm{Sp}_{\mathsf{B}}(x)\)，从而等于 \(\mathrm{Sp}_{\mathsf{A}}(x)\)。由于 B 的特征标都是 Hermite 的 (I, p.~107, 引理 9)，在 \(x\) 上取值相同的 B 的特征标在 \(x^{*}\) 上也取相同的值，因而在由 \(x\) 生成的幺星子代数 B 上相等。这就证明映射 \(\mathrm{ev}_x\) 是单射。由于 \(\mathsf{X}(\mathsf{B})\) 是紧的且 \(\mathbf{C}\) 是分离的，它诱导出从 \(\mathsf{X}(\mathsf{B})\) 到其像上的一个同胚，由此得引理。

由该引理我们导出星代数的一个同构 \(\varphi_x\colon \mathcal{C}(\mathrm{Sp}_{\mathrm{A}}(x))\to \mathcal{C}(\mathrm{X(B)})\)。它把函数 \(f\in \mathcal{C}(\mathrm{Sp}_{\mathrm{A}}(x))\) 映为函数 \(\chi \mapsto f(\chi (x))\)。映射 \(\mathcal{G}_{\mathrm{B}}^{-1}\circ \varphi_{x}\) 是从星代数 \(\mathcal{C}(\mathrm{Sp}_{\mathrm{A}}(x))\) 到 B 上的一个同构。

定义 5. --- 对合态射 \(f \mapsto (\mathcal{G}_{\mathrm{B}}^{-1} \circ \varphi_x)(f)\)（由 \(\mathcal{C}(\mathrm{Sp}_{\mathrm{A}}(x))\) 到 A）称为 \(x\) 在 A 中的连续函数演算映射。我们把它记作 \(f \mapsto f(x)\)。

注记. --- 映射 \(f \mapsto f(x)\) 是等距的；它的像是由 x 生成的幺星子代数 B，它包含在 x 的双交换子中。

若 \(f\) 在 \(\mathrm{Sp}_{\mathrm{A}}(x)\) 上的限制是形如 \(z \mapsto \mathrm{P}(z, \overline{z})\) 的函数（其中 \(\mathrm{P} \in \mathbf{C}[\mathrm{X}, \mathrm{Y}]\) 是一个多项式），则按通常的代数意义有 \(f(x) = \mathrm{P}(x, x^{*})\)。

例. --- 1) 设存在 \(\lambda \in \mathbf{C}\) 使 \(\mathrm{Sp}_{\mathrm{A}}(x)\) 缩减为 \(\lambda\)。那么我们有 \(x = \lambda \cdot 1\)。事实上，\(\mathrm{Sp}_{\mathrm{A}}(x)\) 上的恒等函数等于 \(\lambda\)，因而它在函数演算映射下的像，即 \(x\)，等于 \(\lambda \cdot 1\)。

\begin{enumerate}
\def\labelenumi{\arabic{enumi})}
\setcounter{enumi}{1}
\item
  为使 \(x\) 是 Hermite 的，必须且只需 \(\mathrm{Sp}_{\mathrm{A}}(x)\) 包含在 \(\mathbf{R}\) 中。事实上，设 \(f\) 是 \(\mathrm{Sp}_{\mathrm{A}}(x)\) 上由 \(f(z) = z - \overline{z}\) 给定的连续映射。那么 \(x\) 是 Hermite 的当且仅当 \(f(x) = 0\)，即当且仅当 \(f\) 为零，亦即当且仅当 \(\mathrm{Sp}_{\mathrm{A}}(x)\) 包含在 \(\mathbf{R}\) 中。
\item
  为使 \(x\) 是酉元，必须且只需它的谱包含在 \(\mathbf{C}\) 的单位圆中。事实上，设 \(f \in \mathcal{C}(\mathrm{Sp}_{\mathrm{A}}(x))\) 是由 \(f(z) = z\overline{z} - 1\) 定义的函数；元素 \(x\) 是酉元当且仅当 \(f(x) = 0\)，即当且仅当 \(f\) 为零。
\end{enumerate}

命题 7. — 映射 \(f \mapsto f(x)\) 是从 \(\mathcal{C}(\mathrm{Sp}_{\mathrm{A}}(x))\) 到 A 中的唯一的幺对合代数态射，使得恒等映射 \(z\)（在 \(\mathrm{Sp}_{\mathrm{A}}(x)\) 上）以 \(x\) 为像。

事实上，由 \(\mathcal{C}(\mathrm{Sp}_{\mathrm{A}}(x))\) 的元素 \(z\) 与 \(\overline{z}\) 生成的 \(\mathcal{C}(\mathrm{Sp}_{\mathrm{A}}(x))\) 的幺子代数在 \(\mathcal{C}(\mathrm{Sp}_{\mathrm{A}}(x))\) 中稠密 (TG, X, p.~40, 推论 1)。由于从 \(\mathcal{C}(\mathrm{Sp}_{\mathrm{A}}(x))\) 到 A 的每个对合代数态射都是连续的 (I, p.~104, 命题 2)，至多存在一个从 \(\mathcal{C}(\mathrm{Sp}_{\mathrm{A}}(x))\) 到 A 的对合代数态射，它把 \(z\) 映为 \(x\)。

下面的推论表明：当 \(f\) 是 \(\mathrm{Sp}_{\mathrm{A}}(x)\) 的一个邻域中全纯函数的限制时，\(f(x)\) 的这一定义与 I, p.~74 第 9 小节中单变量全纯函数演算的定义一致。

推论 1. --- 设 \(f \in \mathcal{O}(\mathrm{Sp}_{\mathrm{A}}(x))\) 是 \(\mathrm{Sp}_{\mathrm{A}}(x)\) 的一个邻域中全纯函数的一个芽，\(\widetilde{f} \in \mathcal{C}(\mathrm{Sp}_{\mathrm{A}}(x))\) 是 \(\mathrm{Sp}_{\mathrm{A}}(x)\) 上的与 \(f\) 相伴随的连续函数。我们有 \(\widetilde{f}(x) = f(x)\)，其中 \(f(x)\) 是由全纯函数演算给出的 A 的元素。

事实上，映射 \(f \mapsto \widetilde{f}(x)\) 是从 \(\mathcal{O}(\mathrm{Sp}_{\mathrm{A}}(x))\) 到 A 的一个连续幺态射，它把 \(\mathrm{Sp}_{\mathrm{A}}(x)\) 的一个邻域中恒等函数的芽映为 \(x\)。于是结果由 I, p.~74 的定理 5 得出。

推论 2. --- 设 \(f \in \mathcal{C}(\mathrm{Sp}_{\mathrm{A}}(x))\)。

\begin{enumerate}
\def\labelenumi{\alph{enumi})}
\tightlist
\item
  我们有
\end{enumerate}

\[
\operatorname{Sp} _ {\mathrm{A}} (f (x)) = f (\operatorname{Sp} _ {\mathrm{A}} (x));
\]

\begin{enumerate}
\def\labelenumi{\alph{enumi})}
\setcounter{enumi}{1}
\tightlist
\item
  对每个 \(g \in \mathcal{C}(\mathrm{Sp}_{\mathrm{A}}(f(x)))\)，有 \((g \circ f)(x) = g(f(x))\)。
\end{enumerate}

由于 \(f(x)\) 属于 A 的全子代数 B，我们有 \(\mathrm{Sp}_{\mathrm{A}}(f(x)) = \mathrm{Sp}_{\mathrm{B}}(f(x))\) (I, p.~106, 命题 5)。同构 \(f \mapsto f(x)\)（从 \(\mathcal{C}(\mathrm{Sp}_{\mathrm{A}}(x))\) 到 B）保持谱；因此我们有 \(\mathrm{Sp}_{\mathrm{B}}(f(x)) = \mathrm{Sp}_{\mathcal{C}(\mathrm{Sp}_{\mathrm{A}}(x))}(f) = f(\mathrm{Sp}_{\mathrm{A}}(x))\) (I, p.~17, 例 3)。这就证明了断言 \(a\))。

映射 \(g \mapsto (g \circ f)(x)\) 是从 \(\mathcal{C}(\mathrm{Sp}_{\mathrm{A}}(f(x)))\) 到 A 的一个幺对合代数态射，它把 \(\mathrm{Sp}_{\mathrm{A}}(f(x))\) 的恒等函数映为 \(f(x)\)。于是依命题 7，\((g \circ f)(x) = g(f(x))\) 对一切 \(g \in \mathcal{C}(\mathrm{Sp}_{\mathrm{A}}(f(x)))\) 成立。

例 4. --- 设 X 是一个局部紧空间，A = \(\mathcal{C}_{b}(X)\) 是 X 上有界连续函数组成的交换幺星代数 (I, p.~102, 例 2)。设 \(g \in A\)；它的谱 S 是 g 的值集 \(g(\mathrm{X})\) 在 C 中的闭包 (I, p.~17, 例 3)。于是 g 的函数演算映射就是映射 \(f \mapsto f \circ g\)（对 \(f \in \mathcal{C}(\mathrm{S})\)）。事实上，这个映射是星代数的一个幺态射，且使 S 的恒等映射以 g 为像。

在 X 是紧的情形，我们有 \(A = \mathcal{C}(X)\) 且 \(S = g(X)\)。

设 \(\pi\colon\mathrm{A}\to\mathrm{A}^{\prime}\) 是幺星代数的一个幺态射。元素 \(\pi(x)\) 在 \(\mathrm{A}^{\prime}\) 中是正规的，且它相对于 \(\mathrm{A}^{\prime}\) 的谱含于 \(\mathrm{Sp}_{\mathrm{A}}(x)\) 中。于是我们有：

命题 8. --- 设 \(f \in \mathcal{C}(\mathrm{Sp}_{\mathrm{A}}(x))\)。仍用 \(f\) 表示 \(f\) 在 \(\mathrm{Sp}_{\mathrm{A}'}(\pi(x))\) 上的限制，我们有等式 \(\pi(f(x)) = f(\pi(x))\)。特别地，对每个 \(\chi \in \mathsf{X}(\mathrm{A})\)，有 \(\chi(f(x)) = f(\chi(x))\)。

设 \(z\) 是 \(\mathrm{Sp}_{\mathrm{A}}(x)\) 的恒等映射。由 \(f \mapsto \pi(f(x))\) 与 \(f \mapsto f(\pi(x))\) 定义的映射都是从 \(\mathcal{C}(\mathrm{Sp}_{\mathrm{A}}(x))\) 到 B 的连续幺对合代数态射，且都把 \(z\) 映为 \(\pi(x)\)。因此这两个态射在 \(\mathcal{C}(\mathrm{Sp}_{\mathrm{A}}(x))\) 的由 \(z\) 生成的幺对合子代数上一致。由于后者在 \(\mathcal{C}(\mathrm{Sp}_{\mathrm{A}}(x))\) 中稠密 (TG, X, p.~40, 推论 1)，这两个态射相等。

推论. --- 设 A 是交换的。对每个 \(f \in \mathcal{C}(\mathrm{Sp}_{\mathrm{A}}(x))\)，有 \(\mathcal{G}_{\mathrm{A}}(f(x)) = f \circ \mathcal{G}_{\mathrm{A}}(x)\)。

只需对 Gelfand 变换 \(\mathcal{G}_{\mathrm{A}}\)（从 A 到 \(\mathcal{C}(\mathsf{X}(\mathrm{A}))\)）应用命题 8，并注意（见上面的例）\(f(\mathcal{G}_{\mathrm{A}}(x)) = f \circ \mathcal{G}_{\mathrm{A}}(x)\)。

\subsection{函数演算的应用}

命题 9. --- 星代数的每个单射态射都是等距的，且特别地有闭像。

设 A 与 A′ 是星代数，\(\pi\colon A\to A'\) 是从 A 到 A' 的一个对合代数态射。

首先设 A 与 A' 都是幺的且 \(\pi\) 是幺态射。我们有 \(\|\pi\| \leqslant 1\)（命题 2）。设存在 A 中的 x 使 \(\|\pi(x)\| < \|x\|\)。令 \(y = x^{*}x\)；它是 A 的一个 Hermite 元。由于 A 与 A' 都是星代数，我们有 \(\|\pi(y)\| = \|\pi(x)\|^{2} < \|x\|^{2} = \|y\|\)，即 \(\varrho(\pi(y)) < \varrho(y)\) (I, p.~104, 引理 7)。特别地，\(\mathrm{Sp}_{\mathrm{A}'}(\pi(y))\) 是 \(\mathrm{Sp}_{\mathrm{A}}(y)\) 的一个闭子集，且异于 \(\mathrm{Sp}_{\mathrm{A}}(y)\)

(I, p.~3 的注记 6 与 I, p.~24 的定理 1)。于是存在非零函数 \(f \in \mathcal{C}(\mathrm{Sp}_{\mathrm{A}}(y))\) 使 \(f| \mathrm{Sp}_{\mathrm{A}'}(\pi(y)) = 0\) (TG, IX, p.~13, 命题 3)。令 \(w = f(y) \in \mathrm{A}\)。我们有 \(w \neq 0\)，因为 \(f \neq 0\)；但 \(\pi(w) = \pi(f(y)) = f(\pi(y)) = 0\)，因为 \(f\) 在 \(\mathrm{Sp}_{\mathrm{A}'}(\pi(y))\) 上为零（命题 8）。故 \(\pi\) 不是单射。

现在处理一般情形。设 \(\widetilde{\mathrm{A}}\) 与 \(\widetilde{\mathrm{A}}'\) 分别是由 A 与 \(\mathrm{A}'\) 添加单位元而得的星代数 (I, p.~106, 定义 4)。存在唯一的幺对合代数态射 \(\widetilde{\pi}\colon \widetilde{\mathrm{A}}\to \widetilde{\mathrm{A}}'\)，它拓展 \(\pi\)。这个态射是单射，因而由前所证是等距的。于是对每个 \(x\in \mathrm{A}\)，有 \(\| \pi (x)\| = \| \widetilde{\pi} (x)\| = \| x\|\)。

引理 11. --- 设 X 是一个完全正则拓扑空间，即可一致化且分离的拓扑空间 (TG, IX, p.~8, 定义 4)，至少含两个点。存在 \(\mathcal{C}(\mathrm{X})\) 中非零的连续函数 f 与 g 使 \(fg = 0\)。

设 \(x \neq y\) 是 X 的两个不同点。设 U 与 V 分别是 x 与 y 的开邻域，使得 U 与 V 不相交。由于 X 是可一致化的，依 TG, IX, p.~7, 定理 2，存在函数 \(f \in \mathcal{C}(X)\) 使 \(f(x) = 1\) 且 \(f|X - U = 0\)。类似地，存在 \(g \in \mathcal{C}(X)\) 使 \(g(y) = 1\) 且 \(g|X - V = 0\)。于是我们有 fg = 0。

命题 10. --- 设 A 是一个幺星代数。假设对 A 的每对可交换元素 \((x,y)\)，条件 xy = 0 蕴含 x = 0 或 y = 0。那么 A = C · 1。

若 A 不等于 C·1，则存在 A 中的 Hermite 元 x 不属于 C·1 (I, p.~96, 引理 2)。设 B 是 A 的由 x 生成的幺星子代数。它是交换的，且同构于 \(\mathcal{C}(\mathrm{Sp}_{\mathrm{A}}(x))\) (I, p.~110, 注记)。由于 x 不是标量，它在 B 中的谱不缩减为单个元素 (I, p.~110, 例 1)。因此存在 \(\mathrm{Sp}_{\mathrm{A}}(x)\) 上非零的连续函数 f 与 g 使 fg = 0（引理 11）。A 的元素 \(f(x)\) 与 \(g(x)\) 非零、可交换，且在 A 中满足 \(f(x)g(x) = 0\)。

命题 11. --- 设 A 是一个幺星代数，a、x 与 y 是 A 的元素。设 x 与 y 是正规元。若 xa = ay，则对 x 的谱与 y 的谱的并上每个连续函数 f，有 \(f(x)a = af(y)\)。特别地，\(f(aa^{*})a = af(a^{*}a)\) 对每个函数 \(f \in \mathcal{C}(\mathrm{Sp}'(a^{*}a))\) 成立。

设 \(\mathrm{S} = \mathrm{Sp}(x) \cup \mathrm{Sp}(y)\)。I, p.~106 的命题 6 蕴含 \(x^{*}a = ay^{*}\)。从而，我们得到 \(f(x)a = af(y)\)，对每个函数 \(f\)（形如 \(z \mapsto \mathrm{P}(z,\overline{z})\)，其中 \(\mathrm{P} \in \mathbf{C}[\mathrm{X},\mathrm{Y}]\) 是一个多项式）。由于由函数 \(f \in \mathcal{C}(\mathrm{S})\)（满足 \(f(x)a = af(y)\)）组成的集合是 \(\mathcal{C}(\mathrm{S})\) 的一个闭子代数，依 TG, X, p.~40, 推论 1，它与 \(\mathcal{C}(\mathrm{S})\) 重合。

第二个断言是第一个断言应用于 Hermite 元 \(x = aa^{*}\) 与 \(y = a^{*}a\) 的推论，其中用到 \(\mathrm{Sp}'(a^{*}a) = \mathrm{Sp}'(aa^{*})\) 这一事实 (I, p.~5, 命题 1)。

\subsection{非幺代数中的函数演算}

设 A 是一个星代数，\(\widetilde{\mathbf{A}}\) 是由 A 添加单位元 \(e\) 而得的幺星代数。用 \(\pi\) 表示 Hermite 特征标 \(x + \lambda e \mapsto \lambda\)（从 \(\widetilde{\mathbf{A}}\) 到 \(\mathbf{C}\) 中）；我们有 \(\operatorname{Ker}(\pi) = \mathbf{A}\)。

设 \(x \in \mathrm{A}\) 是一个正规元。它在 \(\widetilde{\mathrm{A}}\) 中是正规的，且 \(\mathrm{Sp}_{\widetilde{\mathrm{A}}}^{\prime}(x) = \mathrm{Sp}_{\mathrm{A}}^{\prime}(x)\)。用 \(\mathcal{C}'(\mathrm{Sp}_{\mathrm{A}}^{\prime}(x))\) 表示由连续函数 \(f\)（在 \(\mathrm{Sp}_{\mathrm{A}}^{\prime}(x)\) 上、满足 \(f(0) = 0\)）组成的星代数。

设 \(f \in \mathcal{C}(\mathrm{Sp}_{\mathrm{A}}^{\prime}(x))\)。由于 \(\pi(f(x)) = f(\pi(x))\) (I, p.~112, 命题 8)，我们有 \(f(x) \in \mathrm{A}\) 当且仅当 \(f(0) = 0\)。映射 \(f \mapsto f(x)\) 定义了从星代数 \(\mathcal{C}'(\mathrm{Sp}_{\mathrm{A}}^{\prime}(x))\) 到 A 的一个对合代数态射，恒等映射 \(z\)（在 \(\mathrm{Sp}_{\mathrm{A}}^{\prime}(x)\) 上）在其下的像是 \(x\)。这个态射是等距的，它的像是 A 的由 \(x\) 生成的星子代数。

命题 12. --- 映射 \(f \mapsto f(x)\) 是从星代数 \(\mathcal{C}'(\mathrm{Sp}_{\mathrm{A}}'(x))\) 到 A 的唯一的对合代数态射，使得恒等映射 \(z\)（在 \(\mathrm{Sp}_{\mathrm{A}}'(x)\) 上）以 \(x\) 为像。

元素 \(z\) 与 \(\overline{z}\)（\(\mathcal{C}'(\mathrm{Sp}_{\mathrm{A}}'(x))\) 的）生成 \(\mathcal{C}'(\mathrm{Sp}_{\mathrm{A}}'(x))\) 的一个稠密子代数（参见 TG, X, p.~40, 推论 2）。由于从星代数 \(\mathcal{C}'(\mathrm{Sp}_{\mathrm{A}}'(x))\) 到星代数 A 的每个对合代数态射都是连续的 (I, p.~104, 命题 2)，结果成立。

上一小节关于函数演算的结果可以推广到一般情形。我们只叙述它们，而把证明留待读者作相应修改后补全。

命题 13. --- 我们有下列性质：

\begin{enumerate}
\def\labelenumi{\alph{enumi})}
\item
  对每个 \(f \in \mathcal{C}'(\mathrm{Sp}_{\mathrm{A}}'(x))\)，有 \(\mathrm{Sp}_{\mathrm{A}}'(f(x)) = f(\mathrm{Sp}_{\mathrm{A}}'(x))\)；
\item
  对每个 \(f \in \mathcal{C}'(\mathrm{Sp}_{\mathrm{A}}'(x))\) 及每个 \(g \in \mathcal{C}'(\mathrm{Sp}_{\mathrm{A}}'(f(x)))\)，有 \((g \circ f)(x) = g(f(x))\)；
\item
  设 A' 是一个星代数，\(\pi\) 是从 A 到 A' 的一个态射；那么 \(\pi(x)\) 在 A' 中是正规的，有 \(\mathrm{Sp}_{\mathrm{A}'}^{\prime}(\pi(x)) \subset \mathrm{Sp}_{\mathrm{A}}^{\prime}(x)\)，且 \(\pi(f(x)) = f(\pi(x))\) 对每个 \(f \in \mathcal{C}'(\mathrm{Sp}_{\mathrm{A}}^{\prime}(x))\) 成立；
\item
  若 A 是交换的，且 \(f \in \mathcal{C}'(\mathrm{Sp}_{\mathrm{A}}'(x))\)，则 \(\mathcal{G}_{\mathrm{A}}'(f(x)) = f \circ \mathcal{G}_{\mathrm{A}}'(x)\)。
\end{enumerate}

注记. --- 设 A 是一个幺星代数，\(\widetilde{\mathrm{A}}\) 是由 A 添加单位元 \(e\) 而得的幺星代数。对每个 \(x \in \mathrm{A}\)，有 \(\mathrm{Sp}_{\mathrm{A}}^{\prime}(x) = \mathrm{Sp}_{\mathrm{A}}(x) \cup \{0\}\)。设 \(x\) 是 A 的一个正规元。那么它是 \(\widetilde{\mathrm{A}}\) 的一个正规元，于是 A 中有两个函数演算映射，第一个定义在 \(\mathcal{C}(\mathrm{Sp}_{\mathrm{A}}(x))\) 上，第二个定义在 \(\mathcal{C}'(\mathrm{Sp}_{\mathrm{A}}^{\prime}(x))\) 上。设 \(f' \in \mathcal{C}'(\mathrm{Sp}_{\mathrm{A}}^{\prime}(x))\)；若用 \(f\) 表示它在 \(\mathrm{Sp}_{\mathrm{A}}(x)\) 上的限制，则我们有 \(f'(x) = f(x)\)。

\subsection{星代数中的正元素}

定义 6. --- 设 A 是一个星代数。称 A 的元素 x 是正的，如果它是 Hermite 的且 \(\mathrm{Sp}_{\mathrm{A}}^{\prime}(x) \subset \mathbf{R}_{+}\)。用 \(\mathrm{A}_{+}\) 表示 A 的正元素组成的集合。它是 \(\mathrm{A}_{h}\) 的一个子集。

我们记 \(x \geqslant y\)，如果 \(x - y \in \mathrm{A}_{+}\)。

若星代数 A 是幺的，则它的单位元是正的。

若 B 是 A 的一个星子代数，则有 \(B_{+} = B \cap A_{+}\) (I, p.~106 命题 5 的推论)。

若 \(\pi\colon A\to B\) 是星代数的一个态射，则 \(\pi(A_{+})\subset B_{+}\)。

例. --- 1) 设 X 是一个局部紧空间。在 X 上在无穷远处趋于 0 的连续函数组成的星代数 \(\mathcal{C}_{0}(\mathrm{X})\) 中，相应地在 X 上有界连续函数组成的星代数 \(\mathcal{C}_{b}(\mathrm{X})\) 中，函数 \(f\) 是正元素当且仅当它取实值且 \(f(x) \geqslant 0\) 对一切 \(x \in \mathrm{X}\) 成立（参见 I, p.~17, 例 3）。

\begin{enumerate}
\def\labelenumi{\arabic{enumi})}
\setcounter{enumi}{1}
\item
  设 A 是一个交换星代数。设 \(a\) 是 A 的一个元素。由于 \(\mathrm{Sp}_{\mathrm{A}}^{\prime}(x)\) 是 \(\{0\}\) 与 Gelfand 变换 \(\mathcal{G}(a)\) 的像的并 (I, p.~37, 命题 6)，元素 \(a\) 是正的当且仅当 Gelfand 变换 \(\mathcal{G}(a)\) 是一个正函数。
\item
  设 E 是一个复 Hilbert 空间。星代数 \(\mathcal{L}(\mathrm{E})\) (I, p.~102, 例 1) 的元素 x 是正的当且仅当它是 EVT, V, p.~45, 定义 6 意义下 E 的正自同态 (I, p.~138, 命题 8)。
\end{enumerate}

引理 12. --- 设 A 是一个幺星代数，\(x \in A\) 是一个 Hermite 元。

\begin{enumerate}
\def\labelenumi{\alph{enumi})}
\item
  元素 \(x\) 是正的当且仅当 \(\left\| \|x\| \cdot 1 - x \right\| \leqslant \|x\|\)；
\item
  若 \(\| x\| \leqslant 1\)，则 \(x\) 是正的当且仅当 \(\| 1 - x\| \leqslant 1\)。
\item
  若 x 是正的，则 1 - x 是正的当且仅当 \(\|x\| \leqslant 1\)；
\item
  若 x 是正的且 \(y \in A_{+}\) 与 x 交换，则 xy 是正的。
\end{enumerate}

元素 \(x\) 是 Hermite 的，因而是正规的。通过考虑由 \(x\) 生成的交换星子代数，可化归为代数 A 是交换的情形，即对某个局部紧拓扑空间 X 有 \(A = \mathcal{C}_0(X)\) 的情形 (I, p.~108, 定理 1)。于是前三个断言由上面的例 1 立即得出。同样，为证断言 \(d\))，可以考虑由 \(x\) 与 \(y\) 生成的星子代数，它是交换的。

命题 14. --- 设 A 是一个星代数。集合 \(\mathrm{A}_{+}\) 是实 Banach 空间 \(\mathrm{A}_{h}\) 中的一个闭的尖凸锥 (EVT, II, p.~11)。

设 \(\widetilde{\mathsf{A}}\) 是由 \(\mathsf{A}\) 添加单位元而得的星代数。由于 \(\mathsf{A}_{+} = \mathsf{A} \cap \widetilde{\mathsf{A}}_{+}\)，只需对 \(\widetilde{\mathsf{A}}\) 证明该命题。于是可设 \(\mathsf{A}\) 有单位元。

我们有 \(0 \in \mathrm{A}_{+}\)。对每个 \(\lambda \in \mathbf{R}_{+}^{*}\) 及每个 \(x \in \mathrm{A}\)，有 \(\mathrm{Sp}_{\mathrm{A}}^{\prime}(\lambda x) = \lambda \mathrm{Sp}_{\mathrm{A}}^{\prime}(x)\)，这蕴含 \(\mathrm{A}_{+}\) 是实 Banach 空间 \(\mathrm{A}_{h}\) 中的一个锥。

为证 \(A_{+}\) 是凸的，只需证明若 x 与 y 是正的，则 \(x + y \geqslant 0\) (EVT, II, p.~11, 命题 10)。通过位似，只需证明：若 \(x \geqslant 0\) 与 \(y \geqslant 0\) 还满足 \(\|x\| \leqslant 1\)、\(\|y\| \leqslant 1\)，则元素 \(\frac{1}{2}(x + y)\) 是正的。但我们有

\[
\left\| 1 - \frac {1}{2} (x + y) \right\| \leqslant \frac {1}{2} \| 1 - x \| + \frac {1}{2} \| 1 - y \| \leqslant 1
\]

依引理 12 的断言 \(b)\)，而这同一断言表明 \(\frac{1}{2}(x + y)\) 是正的。

最后，引理 12 的断言 \(a\)) 也蕴含 \(\mathrm{A}_{+}\) 是闭的。

由于 \(A_{+}\) 是 \(A_{h}\) 中的一个以 0 为顶点的锥，它是突出的当且仅当 \(\mathrm{A}_{+}\cap(-\mathrm{A}_{+})\) 缩减为 0。但若 \(x\in\mathrm{A}_{+}\cap(-\mathrm{A}_{+})\)，则有 \(\mathrm{Sp}_{\mathrm{A}}^{\prime}(x)=\{0\}\)，于是 \(\varrho(x)=0\)，且 \(\|x\|=0\)。由于 x 是 Hermite 的 (I, p.~104, 引理 7, (2))，由此得 x=0。

命题 14 表明关系 \(x\geqslant y\) 是 \(\mathrm{A}_h\) 上的一个序关系 (EVT, II, p.~13, 命题 13)。

命题 15. — 设 A 是一个星代数。设 x 是 A 的一个正规元。

\begin{enumerate}
\def\labelenumi{\alph{enumi})}
\item
  设 A 是幺的，f 是从 \(\mathrm{Sp}_{\mathrm{A}}(x)\) 到 C 中的一个连续函数。为使 \(f(x)\) 是正的，必须且只需 f 的像包含在 \(R_{+}\) 中；
\item
  设 f 是从 \(\mathrm{Sp}_{\mathrm{A}}^{\prime}(x)\) 到 C 中的一个连续函数，满足 \(f(0)=0\)。为使 \(f(x)\) 是 A 的一个正元素，必须且只需 f 的像包含在 \(R_{+}\) 中。
\end{enumerate}

断言 \(a\)) 由 I, p.~111 推论 2 的断言 \(a\)) 得出，断言 \(b\)) 由 I, p.~114 的命题 13 得出。

设 \(x\) 是星代数 A 的一个 Hermite 元。它的谱包含在 \(\mathbf{R}\) 中 (I, p.~106, 命题 4)。考虑从 \(\mathrm{Sp}_{\mathrm{A}}^{\prime}(x)\) 到 \(\mathbf{R}\) 中的由下式定义的连续函数

\[
f _ {1} \colon t \mapsto \sup (t, 0), \quad f _ {2} \colon t \mapsto \sup (- t, 0), \quad f _ {3} \colon t \mapsto | t |.
\]

我们记

\[
x ^ {+} = f _ {1} (x), \quad x ^ {-} = f _ {2} (x), \quad | x | = f _ {3} (x).\tag{3}
\]

由于函数 \(f_1, f_2, f_3\) 取实值、是正的且在 0 处为零，元素 \(x^+, x^-\) 与 \(|x|\) 是 A 的正元素（命题 15, a)），且属于 A 的由 \(x\) 生成的星子代数。

我们有 \(f_{1}(t) - f_{2}(t) = t\) 对一切 \(t \in \mathbf{R}\) 成立，且有关系式 \(f_{1} + f_{2} = f_{3}\) 与 \(f_{1}f_{2} = 0\)。由此得

\[
x = x ^ {+} - x ^ {-}, \quad | x | = x ^ {+} + x ^ {-}, \quad x ^ {+} x ^ {-} = x ^ {-} x ^ {+} = 0.\tag{4}
\]

由于函数演算映射是等距的，我们有

\[
\| | x | \| = \| x \|, \quad \| x ^ {+} \| \leqslant \| x \|, \quad \| x ^ {-} \| \leqslant \| x \|.
\]

设 \(x\) 是 A 的一个正元素。它是 Hermite 的，因而是正规的。设 \(\alpha \in \mathbf{R}_{+}^{*}\)，并设 \(g\) 是 \(\mathrm{Sp}_{\mathrm{A}}^{\prime}(x)\) 上函数 \(t \mapsto t^{\alpha}\) 的限制；我们记 \(x^{\alpha} = g(x)\)。这是 A 的由 \(x\) 生成的星子代数中的一个正元素。设 \(\alpha\) 与 \(\beta\) 属于 \(\mathbf{R}_{+}^{*}\)。由于函数演算映射是一个代数同态，并依 I, p.~114 的命题 13，我们有

\[
x ^ {\alpha} x ^ {\beta} = x ^ {\alpha + \beta} \quad (x ^ {\alpha}) ^ {\beta} = x ^ {\alpha \beta}.\tag{5}
\]

我们也记作 \(\sqrt{x} = x^{1/2}\)。

命题 16. — 设 x 是 A 的一个正元素。设 \(\alpha \in \mathbf{R}_{+}^{*}\)。那么 \(x^{1/\alpha}\) 是 A 中满足 \(y^{\alpha} = x\) 的唯一的正元素 y。

我们在上面已看到 \(x^{1/\alpha}\) 满足所要求的性质。反之，设 y 是 A 中满足 \(y^{\alpha} = x\) 的一个正元素。依公式 (5)，我们有 \(y = (y^{\alpha})^{1/\alpha} = x^{1/\alpha}\)，这就是要证的。

引理 13. — 设 A 是一个幺星代数。A 的每个元素都是酉元之和。

设 \(x\) 是 A 的一个 Hermite 元。首先设 \(\| x \| \leqslant 2\)。依引理 12, c)，我们有 \(1 - \frac{1}{4} x^2 \in \mathrm{A}_+\)。令 \(y = \frac{1}{2} x + i\sqrt{1 - \frac{1}{4} x^2}\)。我们有 \(y^* = \frac{1}{2} x - i\sqrt{1 - \frac{1}{4} x^2}\)，于是 \(yy^* = 1\)，且 \(x = y + y^*\) 是两个酉元之和。在一般情形，取整数 \(k\) 使 \(\| \frac{1}{k} x \| \leqslant 2\)；于是元素 \(x\) 是 \(2k\) 个酉元之和。依 I, p.~96 的引理 2，引理得证。

定理 2. — 设 A 是一个星代数。A 的元素 x 是正的当且仅当存在 \(y \in A\) 使 \(x = y^{*}y\)。

设 \(x\) 是正的。令 \(y = x^{1/2}\)；它是 A 的一个 Hermite 元，且我们有 \(y^*y = y^2 = x\)。

反之，设 y 是 A 的一个元素，并令 \(x = y^{*}y\)。这是 A 的一个 Hermite 元。我们来证明 x 是正的。为此，记 \(z = x^{-}\) 并令 w = yz。于是我们有

\[
w ^ {*} w = z ^ {*} y ^ {*} y z = z x z = z (x ^ {+} - z) z = - z ^ {3}.
\]

由于 \(z \geqslant 0\)，有 \(z^3 \geqslant 0\)，于是我们推出 \(\mathrm{Sp}_{\mathrm{A}}'(w^{*}w) \subset \mathbf{R}_{-}\)。再写 \(w = w_1 + iw_2\)，其中 \(w_1\) 与 \(w_2\) 是 Hermite 的 (I, p.~96, 引理 2)。元素 \(w_1^2\) 与 \(w_2^2\) 是正的。我们有 \(ww^{*} + w^{*}w = 2w_{1}^{2} + 2w_{2}^{2}\)，于是命题 14 表明 \(ww^{*} = 2w_{1}^{2} + 2w_{2}^{2} + (-w^{*}w)\) 是正的。由于 \(\mathrm{Sp}_{\mathrm{A}}'(ww^{*}) = \mathrm{Sp}_{\mathrm{A}}'(w^{*}w)\) (I, p.~5, 命题 1)，而后者包含在 \(\mathbf{R}_{-}\) 中，我们得出结论 \(\mathrm{Sp}_{\mathrm{A}}'(w^{*}w) = \{0\}\)。由于 \(w^{*}w\) 是 Hermite 的，这蕴含 (I, p.~108, 推论 1) \(\|w^{*}w\| = \varrho(w^{*}w) = 0\)，从而 \(z^{3} = 0\)。由于 z 是 Hermite 的，我们有 z = 0。于是 \(x = x^{+}\) 是正的。

注记. --- 设 A 是一个星代数，\(x \in A\)。元素 \(x^{*}x\) 是正的；因此我们令 \(|x| = (x^{*}x)^{1/2}\)。我们有 \(\|x\|^{2} = \|x^{*}x\| = |||x|^{2}\| = |||x||^{2}\)，从而 \(\|x\| = |||x|\|\)。

当 x 是正规元时，我们还有 \(|x| = f(x)\)，其中 f 是从 C 到 C 的在 0 处为零、由 \(f(z) = |z|\) 给定的映射。特别地，当 x 是 Hermite 元时，\(|x|\) 与公式 (3) 定义的元素重合。

进一步设 A 是幺的且 x 是可逆的。那么 \(|x|\) 也是可逆的，元素 \(u = x|x|^{-1}\) 是酉元，且我们有 \(x = u|x|\)（极分解；关于 Hilbert 空间自同态的情形另见 I, p.~139, 第 8 小节）。

引理 14. --- 设 A 是一个星代数，x 与 y 是 A 的 Hermite 元。

\begin{enumerate}
\def\labelenumi{\alph{enumi})}
\item
  若 \(x \leqslant y\)，则对 A 的每个元素 \(w\)，有 \(w^{*}xw \leqslant w^{*}yw\)。特别地，若 \(y \geqslant 0\)，则有 \(w^{*}yw \geqslant 0\)；
\item
  设 A 是幺的。若 \(0 \leqslant x \leqslant y\) 且 \(x\) 可逆，则 \(y\) 可逆且 \(y^{-1} \leqslant x^{-1}\)；
\item
  若 \(0 \leqslant x \leqslant y\)，则 \(\|x\| \leqslant \|y\|\)。
\end{enumerate}

令 \(u = (y - x)^{1/2}\)。我们有 \(w^{*}yw - w^{*}xw = w^{*}u^{2}w = (uw)^{*}(uw)\)，于是断言 a) 由定理 2 得出。

我们来证明断言 \(b\))。首先设 \(x = 1\)。设 B 是由 \(y\) 生成的幺星子代数。依 Gelfand 同构，\(y\) 对应于一个 \(\geqslant 1\) 的连续函数，定义在紧空间 \(\mathsf{X}(\mathsf{B})\) 上。于是这个函数是可逆的，且它的逆 \(\leqslant 1\)。这蕴含 \(y\) 可逆且 \(y^{-1} \leqslant 1 = x^{-1}\)。在一般情形，我们注意到 \(0 \leqslant 1 \leqslant x^{-1/2}yx^{-1/2}\)（依 \(a\))），故前面的情形蕴含 \(z = x^{-1/2}yx^{-1/2}\) 可逆且 \(z^{-1} \leqslant 1\)。于是 \(y\) 可逆且 \(y^{-1} \leqslant x^{-1}\)，这仍依 \(a\))。

为证断言 \(c\))，可设 A 是幺的 (I, p.~105, 命题 3)。首先设 \(y\) 可逆。令 \(b = \sqrt{y}\)。依 \(a\))，条件 \(0 \leqslant x \leqslant y\) 蕴含 \(0 \leqslant b^{-1}xb^{-1} \leqslant b^{-1}yb^{-1} = 1\)，于是 \(\|b^{-1}xb^{-1}\| \leqslant 1\)，依 I, p.~116 的引理 12 \(c\))。于是我们有

\[
\| x \| = \| b (b ^ {- 1} x b ^ {- 1}) b \| \leqslant \| b \| \| b ^ {- 1} x b ^ {- 1} \| \| b \| \leqslant \| b \| ^ {2} = \| b ^ {2} \| = \| y \|.
\]

在一般情形，对每个实数 \(\varepsilon > 0\)，元素 \(y + \varepsilon\) 可逆且 \(0 \leqslant x \leqslant y + \varepsilon\)。于是由前所证，\(\|x\| \leqslant \|y + \varepsilon\|\) 对每个实数 \(\varepsilon > 0\) 成立，由此得结果。

注. — 一般地，若 x 与 y 是星代数 A 的正元素，条件 \(0 \leqslant x \leqslant y\) 并不蕴含 \(x^{2} \leqslant y^{2}\)（参见 I, p. 182, 习题 15）。

\subsection{星代数中的近似单位元}

定义 7. --- 设 A 是一个赋范代数。A 的一个近似单位元是 A 的单位球上的一个滤子基 F，使得对 A 中每个 x，A 上的滤子基 xF 与 Fx 收敛于 x，换言之：

\[
\lim _ {f, \mathfrak {F}} f x = \lim _ {f, \mathfrak {F}} x f = x.
\]

若 A 是一个星代数，则称近似单位元 F 是递增的，如果 F 是 \(A_{+}\) 上的一个滤子基。

设 A 是一个星代数。我们用 \(A_{+}^{\leqslant1}\)（相应地，\(A_{+}^{<1}\)）表示 A 中范数 \(\leqslant 1\)（相应地，范数 \(< 1\)）的正元素组成的集合；它们是 A 的谱包含在 \([0,1]\)（相应地，\([0,1)\)）中的 Hermite 元。

命题 17. --- 设 A 是一个星代数，\(\mathfrak{F}\) 是 \(\mathrm{A}_{+}^{\leqslant 1}\) 上的一个滤子基。为使 \(\mathfrak{F}\) 是 A 的一个递增近似单位元，必须且只需有

\[
\lim _ {f, \mathfrak {F}} f x = x
\]

对 A 的每个正元素 x 成立。

这个条件显然是必要的；我们来证明它是充分的。设 \(\widetilde{\mathrm{A}}\) 是由 A 添加单位元而得的星代数。设 \(x\) 是 A 的一个元素，\(f \in \mathrm{A}_{+}^{\leqslant 1}\)。我们有

\[
\begin{array}{c} \| f x - x \| ^ {2} = \| (f x - x) (f x - x) ^ {*} \| = \| (f - 1) x x ^ {*} (f - 1) \| \\ \leqslant \| (f - 1) x x ^ {*} \| \end{array}
\]

因为 \(\| f - 1\| \leqslant 1\)（I, p.~116 的引理 12, \(c\))）。于是我们有

\[
\limsup _ {f, \mathfrak {F}} \| f x - x \| ^ {2} \leqslant \limsup _ {f, \mathfrak {F}} \| f x x ^ {*} - x x ^ {*} \|.
\]

由于 \(xx^{*}\) 是正的 (I, p.~118, 定理 2)，假设蕴含

\[
\limsup _ {f, \mathfrak {F}} \| f x - x \| ^ {2} \leqslant \| x x ^ {*} - x x ^ {*} \| = 0,
\]

从而

\[
\lim _ {f, \mathfrak {F}} f x = x.
\]

由于 A 的对合是连续的且 \(\mathfrak{F}\) 是 \(A_{h}\) 上的一个滤子基，由此得

\[
\lim _ {f, \mathfrak {F}} x f = \lim _ {f, \mathfrak {F}} (f ^ {*} x ^ {*}) ^ {*} = \lim _ {f, \mathfrak {F}} (f x ^ {*}) ^ {*} = (x ^ {*}) ^ {*} = x.
\]

命题由此得证。

命题 18. --- 设 A 是一个星代数。有序集 \(A_{+}^{<1}\) 是向右定向的 (E, III, p.~12, 定义 7)，且它的截口滤子 (TG, I, p.~38, 例 2) 是 A 的一个递增近似单位元。

设 \(\widetilde{\mathsf{A}}\) 是由 A 添加记作 1 的单位元而得的星代数 (I, p.~106, 定义 4)。

设 \(g\) 是从 \([0,1)\) 到 \(\mathbf{R}_+\) 中的由 \(g(t) = t(1 - t)^{-1}\) 定义的函数。它是一个连续递增双射，其逆双射由 \(t \mapsto 1 - (1 + t)^{-1}\) 给出。

我们来证明有序集 \(A_{+}^{<1}\) 是向右定向的。设 x 与 y 是 \(A_{+}^{<1}\) 的元素。由于 \(g(0)=0\) 且 \(\mathrm{Sp}_{\mathrm{A}}^{\prime}(x)\) 与 \(\mathrm{Sp}_{\mathrm{A}}^{\prime}(y)\) 都包含在 \([0,1)\) 中，元素 \(g(x)\) 与 \(g(y)\) 有定义；它们是正的，于是 \(g(x)+g(y)\geqslant0\)。因此我们可以作 A 的元素 \(z=g^{-1}(g(x)+g(y))\)。我们有 \(\mathrm{Sp}_{\mathrm{A}}^{\prime}(z)\subset[0,1)\)，从而 \(z\in A_{+}^{<1}\)。

我们有 \(0 \leqslant g(x) \leqslant g(z)\)，于是 \(1 \leqslant 1 + g(x) \leqslant 1 + g(z)\)。I, p.~119 的引理 14 的断言 \(b\)) 蕴含 \(1 + g(x)\) 与 \(1 + g(z)\) 可逆，且 \((1 + g(z))^{-1} \leqslant (1 + g(x))^{-1}\)。从而 \(z = 1 - (1 + g(z))^{-1} \geqslant 1 - (1 + g(x))^{-1} = x\)。类似地，我们有 \(z \geqslant y\)。因此，\(z\) 优超 \(x\) 与 \(y\)（在 \(\mathrm{A}_{+}^{<1}\) 中）。于是有序集 \(\mathrm{A}_{+}^{<1}\) 是向右定向的。用 \(\mathfrak{F}\) 表示它的截口滤子

设 \(x\) 是 A 的一个正元素。对每个整数 \(n \geqslant 1\)，令 \(e_n = g^{-1}(nx)\)；我们有 \(e_n \in \mathrm{A}_+^{<1}\)。设 \(h_n\) 是 \(\mathbf{R}_+\) 上对一切 \(t \in \mathbf{R}_+\) 由 \(h_n(t) = t^2(1 - g^{-1}(nt)) = t^2 / (1 + nt)\) 定义的连续函数。我们有 \(|h_n(t)| \leqslant t / n\) 对一切 \(t \geqslant 0\) 成立，从而 \(\| x(1 - e_n)x \| = \| h_n(x) \| \leqslant \| x \| / n\)。特别地，\(x(1 - e_n)x\) 当 \(n\) 趋于无穷时趋于 0。

设 \(\varepsilon > 0\) 是一个实数。取整数 \(n\) 使 \(\|x(1 - e_n)x\| < \varepsilon\)。于是对每个 \(f \in \mathrm{A}_+^{<1}\)，只要 \(f \geqslant e_n\)，我们有

\[
\begin{array}{r l} & {\| x - f x \| ^ {2} = \| (1 - f) x \| ^ {2} = \| ((1 - f) x) ^ {*} (1 - f) x \| = \| x ^ {*} (1 - f) ^ {2} x \|} \\ & {\qquad \qquad \qquad \qquad \qquad \qquad \qquad \qquad \qquad \qquad = \| x (1 - f) ^ {2} x \|.} \end{array}
\]

此外，由于 \(0 \leqslant f \leqslant 1\)，由此得 \((1-f)-(1-f)^{2} = (1-f)f \geqslant 0\)（I, p.~116 的引理 12, d)），从而 \((1-f)^{2} \leqslant 1-f\)。由于 \(1-f \leqslant 1-e_{n}\)，由此得 \(0 \leqslant (1-f)^{2} \leqslant 1-e_{n}\)。依 I, p.~119 的引理 14, a) 与 c)，由此得

\[
\left\| x - f x \right\| ^ {2} = \left\| x (1 - f) ^ {2} x \right\| \leqslant \left\| x (1 - e _ {n}) x \right\| <   \varepsilon .
\]

于是 \(\lim_{f,\mathfrak{F}}fx = x\) 对一切 \(x \in A_{+}\) 成立。因此依命题 17，滤子 \(\mathfrak{F}\) 是 A 的一个近似单位元。

\subsection{对闭双边理想的商}

引理 15. --- 设 A 是一个星代数，I 是 A 的一个闭双边理想。那么 I 是自伴的。

令 \(J = I \cap I^*\)。集合 \(J\) 是 \(A\) 的一个自伴双边理想，它包含 \(I^*I\)。特别地，\(J\) 是一个星代数。设 \(\mathfrak{F}\) 是 \(J\) 的一个递增近似单位元 (I, p.~121, 命题 18)。对一切 \(x \in I\) 及每个 \(f \in J_{+}^{\leqslant 1}\)，有

\[
\begin{array}{c} \| x f - x \| ^ {2} = \| (x f - x) ^ {*} (x f - x) \| = \| f (x ^ {*} x f - x ^ {*} x) - (x ^ {*} x f - x ^ {*} x) \| \\ \leqslant 2 \| x ^ {*} x f - x ^ {*} x \|. \end{array}
\]

由于 \(x^{*}x\in \mathrm{J}\)，我们因此有

\[
\lim _ {f, \mathfrak {F}} \| x f - x \| ^ {2} = 0.
\]

由于 \(xf \in \mathrm{J}\) 对一切 \(f \in \mathrm{J}_{+}^{\leqslant 1}\) 成立，且 \(\mathrm{J}\) 是闭的，这蕴含 \(x \in \mathrm{J}\)。因此 \(\mathrm{I} = \mathrm{J}\)，即理想 \(\mathrm{I}\) 是自伴的。

命题 19. --- 设 A 是一个星代数，I 是 A 的一个闭双边理想。那么商对合 Banach 代数 A/I 是一个星代数。

理想 I 是自伴的（引理 15）。我们考虑由 A 添加单位元而得的星代数 \(\widetilde{\mathsf{A}}\) (I, p.~106, 定义 4)。在这个代数中，集合 I 是一个闭自伴双边理想，且 A/I 恒等同于 \(\widetilde{\mathrm{A}}/\mathrm{I}\) 的一个闭对合子代数。于是可设 A 是幺的。

Banach 代数 A/I 是对合的。设 \(\pi\colon\mathrm{A}\to\mathrm{A}/\mathrm{I}\) 是典范投影。自伴双边理想 I 是 A 的一个星子代数。设 \(\mathfrak{F}\) 是 I 的一个递增近似单位元 (I, p.~121, 命题 18)。我们首先证明对每个 \(x\in\mathrm{A}\)，有

\[
\| \pi (x) \| _ {\mathrm{A/I}} = \lim _ {f, \mathfrak {F}} \| x - x f \|.\tag{6}
\]

一方面，由于 \(xf \in \mathrm{I}\) 对一切 \(f \in \mathrm{I}_{+}^{\leqslant 1}\) 成立，有

\[
\| \pi (x) \| _ {\mathrm{A} / \mathrm{I}} = \inf _ {a \in \mathrm{I}} \| x - a \| \leqslant \operatorname * {l i m i n f} _ {f, \mathfrak {F}} \| x - x f \|.
\]

另一方面，对每个 \(a \in \mathrm{I}\)，有

\[
\begin{array}{c} \| x - x f \| \leqslant \| (x - a) - (x - a) f \| + \| a - a f \| \\ = \| (x - a) (1 - f) \| + \| a - a f \| \end{array}
\]

因而，由于 \(\|1 - f\| \leqslant 1\) (I, p.~116, 引理 12) 且 \(a \in I\)，我们导出

\[
\operatorname * {l i m s u p} _ {f, \mathfrak {F}} \| x - x f \| \leqslant \| x - a \|.
\]

于是，得出 \(\limsup_{f,\mathfrak{F}} \|x - xf\| \leqslant \|\pi(x)\|_{\mathrm{A/I}}\)，因为 \(a\) 在 I 中任意。公式 (6) 因此得证。

现设 \(x\) 是 A 的一个元素。依公式 (6)，我们有

\[
\begin{array}{c} \| \pi (x) \| _ {\mathrm{A/I}} ^ {2} = \lim _ {f, \mathfrak {F}} \| x - x f \| ^ {2} = \lim _ {f, \mathfrak {F}} \| x (1 - f) \| ^ {2} \\ = \lim _ {f, \mathfrak {F}} \| (1 - f) x ^ {*} x (1 - f) \| \leqslant \lim _ {f, \mathfrak {F}} \| x ^ {*} x (1 - f) \| = \| \pi (x ^ {*} x) \| _ {\mathrm{A/I}}. \end{array}
\]

于是 I, p.~103 的引理 6 蕴含对合 Banach 代数 A/I 是一个星代数。

\subsection{对合 Banach 代数的包络星代数}

引理 16. --- 设 A 是一个对合代数，p 是 A 上的一个半范数。下列条件等价：

\begin{enumerate}
\def\labelenumi{(\roman{enumi})}
\item
  我们有 \(p(xy) \leqslant p(x)p(y)\)，\(p(x^{*}) = p(x)\) 且 \(p(x)^{2} = p(x^{*}x)\) 对一切 \(x, y \in \mathrm{A}\) 成立；
\item
  A 中满足 \(p(x) = 0\) 的元素 x 组成的集合 R 是 A 的一个自伴双边理想，且由 p 在 A/R 上诱导的半范数使 A/R 成为一个对合赋范代数，其完备化是一个星代数（C*-代数）；
\item
  存在一个星代数 B 及从 A 到 B 的一个对合代数态射 \(\varphi\)，使得 \(p(x) = \|\varphi(x)\|\) 对一切 \(x\in\mathbf{A}\) 成立。
\end{enumerate}

蕴含关系 (i) \(\Longrightarrow\) (ii) \(\Longrightarrow\) (iii) \(\Longrightarrow\) (i) 都是初等的。

满足引理 16 中条件的半范数将称为对合代数 A 上的星半范数。

设 A 是一个赋范对合代数，S 是 A 上星半范数组成的集合。我们有 \(p(x) \leqslant \|x\|\) 对一切 \(x \in A\) 及每个 \(p \in S\) 成立 (I, p.~104, 命题 2)。映射 \(x \mapsto \|x\|_* = \sup_{p \in S} p(x)\) 是 A 上的一个星半范数。它是 A 上最大的星半范数。

设 R 是由 \(x \in A\) 中满足 \(\|x\|_{*} = 0\) 的元素组成的集合。它是 A 的一个闭双边理想。我们用 Stell(A) 表示 A/R 对于由 \(x \mapsto \|x\|_{*}\) 诱导的范数的完备化星代数（引理 16, (ii)）。从 A 到 Stell(A) 中的典范映射是连续的，它的像在 Stell(A) 中稠密，且它的核等于 R。

定义 8. --- 星代数 Stell(A) 称为赋范对合代数 A 的包络星代数。

若 A 是交换的，则 Stell(A) 是交换的；若 A 是幺的，则 Stell(A) 是幺的。

命题 20. --- 设 A 是一个赋范对合代数，j 是从 A 到 Stell(A) 中的典范态射。对每个星代数 B 及每个从 A 到 B 的对合代数态射 \(\varphi\)，存在唯一的从 Stell(A) 到 B 中的星代数态射 \(\varphi'\) 使 \(\varphi = \varphi' \circ j\)。

我们用 \(x \mapsto \|x\|_*\) 表示 Stell(A) 上的范数。设 R 是 \(j\) 的核。映射 \(x \mapsto \| \varphi(x) \|\) 是 A 上的一个星半范数。于是 \(\| \varphi(x) \| \leqslant \| x \|_*\) 对一切 \(x \in A\) 成立。因而态射 \(\varphi\) 通过取商定义了从 A/R 到 B 的一个连续态射，它由连续性拓展为从 Stell(A) 到 B 的一个态射 \(\varphi'\)，满足 \(\varphi = \varphi' \circ j\)。\(\varphi'\) 的唯一性源于 \(j\) 的像在 Stell(A) 中稠密这一事实。

推论. --- 设 A 是一个交换对合 Banach 代数，j 是从 A 到 Stell(A) 中的典范态射。映射 \(\mathsf{X}(j)\) 是从 \(\mathsf{X}(\mathrm{Stell}(\mathrm{A}))\) 到子空间 \(\mathsf{X}(\mathrm{A})_{h}\)（由 A 的 Hermite 特征标组成的 \(\mathsf{X}(\mathrm{A})\) 的子空间）上的一个同胚。

A 的 Hermite 特征标就是从 A 到星代数 C 的对合代数态射。于是命题 20 蕴含 \(\mathsf{X}(j)\) 是从 \(\mathsf{X}(\mathrm{Stell}(\mathsf{A}))\) 到 \(\mathsf{X}(\mathsf{A})_h\) 上的一个双射。由于 \(\mathsf{X}(j)\) 是到其像上的一个同胚（参见 I, p.~10），推论得证。

我们把 \(\mathsf{X}(\mathrm{Stell}(\mathbf{A}))\) 与 \(\mathsf{X}(\mathbf{A})_h\) 通过映射 \(\mathsf{X}(j)\) 恒等同。对每个 \(x\in \mathbf{A}\)，映射 \(\mathcal{G}_{\mathrm{Stell(A)}}(j(x))\) 在 \(\mathsf{X}(\mathbf{A})_h\) 上不是别的，正是 \(\mathcal{G}_{\mathbf{A}}(x)\)。

命题 21. --- 设 A 是一个对合 Banach 代数，j 是从 A 到 Stell(A) 中的典范态射。A 的根式包含在 j 的核中。

设 \(x\) 是 A 的根式的一个元素。那么 \(x^{*}x\) 属于 A 的根式，因此 \(\mathrm{Sp}_{\mathrm{A}}^{\prime}(x^{*}x) = \{0\}\) (I, p.~5, 注记 3)。由于 \(\mathrm{Sp}_{\mathrm{Stell(A)}}^{\prime}(j(x^{*}x)) \subset \mathrm{Sp}_{\mathrm{A}}^{\prime}(x^{*}x)\)，我们于是有 \(\mathrm{Sp}_{\mathrm{Stell(A)}}^{\prime}(j(x)^{*}j(x)) = \{0\}\)，从而 \(\| j(x)\| ^2 = \| j(x)^{*}j(x)\| = \varrho (j(x)^{*}j(x)) = 0\) (I, p.~104, 公式 (2))，因此 \(j(x) = 0\)。
